% concatenated from main.tex
\documentclass[12pt]{amsart}
\usepackage{float}
\raggedbottom

% Package Being Used:

\usepackage{amsmath}
\usepackage{amssymb}
\usepackage{esvect}
\usepackage{verbatim}
\usepackage{bm}
\usepackage{graphicx}
\usepackage{psfrag}
\usepackage{color}
\usepackage[dvipsnames]{xcolor}
\usepackage{yfonts}
\usepackage{hyperref}
\hypersetup{colorlinks=true, linkcolor=blue, citecolor=magenta, urlcolor=wine}
\usepackage{url}
\usepackage{algpseudocode}
\usepackage{fancyhdr}
\usepackage{mathtools}
\usepackage{tikz-cd}
\usepackage{stmaryrd}
\usepackage{calrsfs}
\usepackage{quiver}

\usepackage{adjustbox}


% Paper Format and Geometry:

\voffset=-1.4mm
\oddsidemargin=14pt
\evensidemargin=14pt
\topmargin=26pt
\headheight=9pt     
%\textheight=576pt
\textwidth=441pt %441
\parskip=6pt plus 0pt
%\renewcommand{\baselinestretch}{1.5} 

% Head Labels:

\pagestyle{fancy}
\fancyhf{}
\renewcommand{\headrulewidth}{0pt}
\renewcommand{\footrulewidth}{0pt}
\fancyhead[CE]{\fontsize{9}{11}\selectfont SOPHIE ZHU}
\fancyhead[CO]{\fontsize{9}{11}\selectfont RANGES OF TWISTED DIVISOR FUNCTIONS ON NUMBER FIELDS}
\fancyhead[LE,RO]{\thepage}
\newcommand{\newthisweek}[1]{{\color{blue} \sf $\spadesuit$ NEW THIS WEEK: [#1]}}
\newcommand\numberthis{\addtocounter{equation}{1}\tag{\theequation}}
\newcommand\scalemath[2]{\scalebox{#1}{\mbox{\ensuremath{\displaystyle #2}}}}


\setlength{\headheight}{9pt}

% Theorems-like Format and Numbering:

\newtheorem*{maintheorem*}{Main Theorem}
\newtheorem{theorem}{Theorem}[section]
\newtheorem*{theorem*}{Main Theorem}
\newtheorem{prob}[theorem]{Problem}
\newtheorem{question}[theorem]{Question}
\newtheorem{prop}[theorem]{Proposition}
\newtheorem{conj}[theorem]{Conjecture}
\newtheorem{claim}[theorem]{Claim}
\newtheorem{corollary}[theorem]{Corollary}

\newtheorem{lemma}[theorem]{Lemma}
\newtheorem{cor}[theorem]{Corollary}
\newtheorem{exercise}{Exercise}[section]
\newtheorem*{sol}{Solution}

\theoremstyle{definition}
\newtheorem{definition}[theorem]{Definition}
\newtheorem{remark}[theorem]{Remark}
\newtheorem{example}[theorem]{Example}
\numberwithin{equation}{section}

% Personalized Commands:

% New, categorical commands
\newcommand{\CC}{\mathcal C}

\newcommand{\DD}{\mathcal{D}}
\newcommand{\Ob}{\text{Ob}}
\newcommand{\id}{\text{id}}
\newcommand{\set}{\mathsf{Set}}
\newcommand{\hh}{\text{Hom}}
\newcommand{\one}{\mathbb{1}}
\newcommand{\ev}{\text{ev}}
\newcommand{\coev}{\text{coev}}
\newcommand{\vs}{\textsf{Vec}}
\newcommand{\fp}{\textsf{Fpdim}}
 \newcommand{\End}{\textsf{End}}
\newcommand{\Aut}{\textsf{Aut}}
\newcommand{\fa}{\mathfrak{a}}
\newcommand{\bb}{\mathfrak{b}}
\newcommand{\dd}{\mathfrak{d}}

% General commands
\newcommand{\Gal}{\text{Gal}}
\newcommand{\ra}{\rightarrow}
\newcommand{\res}{\text{res}}
\newcommand{\C}{\mathbb{C}}
\newcommand{\ff}{\mathbb{F}}
\newcommand{\pp}{\mathbb{P}}
\newcommand{\p}{\mathfrak{p}}
\newcommand{\QQ}{\mathfrak{Q}}
\newcommand{\q}{\mathfrak{q}}
\newcommand{\Q}{\mathbb{Q}}
\newcommand{\mca}{\mathcal{A}}
\newcommand{\mce}{\mathcal{E}}
\newcommand{\mcr}{\mathcal{R}}
\newcommand{\aaa}{\mathbf{A}}
\newcommand{\uu}{\mathcal{U}}
\newcommand{\Z}{\mathbb{Z}}
\newcommand{\N}{\mathbb{N}}
\newcommand{\R}{\mathbb{R}}
\newcommand{\PP}{\mathbb{P}}
\newcommand{\Frob}{\operatorname{Frob}}

\newcommand{\lcm}{\text{lcm}}
\newcommand{\dash}{\text{-}}
\providecommand\ldb{\llbracket}
\providecommand\rdb{\rrbracket}
\newcommand\pval{\mathsf{v}_p}
\newcommand{\gp}{\text{gp}}
\newcommand{\qf}{\text{qf}}
\newcommand{\supp}{\textsf{supp}}
\newcommand{\ii}{\mathcal{Irr}}
\newcommand{\zzz}{\text{Int}(\zz)}
\newcommand{\rr}{\mathbb{R}}
\newcommand{\nn}{\mathbb{N}}

\colorlet{darkgreen}{green!35!black}

% \keywords{atomicity, atomic monoid, commutative semiring, Laurent semiring, Laurent polynomial, ACCP, bounded factorization monoid, BFM, finite factorization monoid, FFM, half-factorial monoid, HFM, length-factorial monoid, LFM, elasticity}

% \subjclass[2010]{Primary: 20M13; Secondary: 20M14, 16Y60}

\begin{document}
	
	\mbox{}
	\title{Connected components of the ranges of twisted divisor functions on number fields}
	\author{Sophie Zhu}
	\address{Harvard College}
	\email{sophiezhu@college.harvard.edu}
% \author{}
\date{}

\begin{abstract} 
    Let $r\in\mathbb{C}$, let $K$ be a finite extension of $\mathbb{Q}$, let $I_K$ be the monoid of integral ideals in the ring of integers $\mathcal{O}_K$ of $K$, and let $\chi$ be a Dirichlet character. Then define the twisted ideal divisor function $\sigma_{r, K, \chi} : I_K \rightarrow \mathbb{C}$ by $$\sigma_{r,K,\chi}(I) = \sum_{J \mid I} N(J)^{-r}\chi(N(J)),$$ where $N$ denotes the ideal norm. For real $r>1,$ we study the number of connected components $C_{r, K, \chi}$ of the closure $\overline{\sigma_{r,K,\chi}(I_K)}$, writing $C_{r,K}$ when $\chi$ is the principal character modulo 1. We prove that $C_{r,K,\chi}$ is finite when $\chi$ is real-valued. When $K = \mathbb{Q}$, we show that for fixed $r > 1,$ every sufficiently large positive integer is realized as $C_{r,\Q,\chi},$ and if $r$ is sufficiently large, then every positive integer is realized as $\chi$ varies. For finite Galois extensions $K$ over $\mathbb{Q}$, we exhibit new exponential lower bounds for $C_{r,K},$ and we prove that for every fixed integer $s \geq 2$, the values $C_{r,K}$ are unbounded as $K$ ranges over degree-$s$ extensions of $\mathbb{Q}$. 
\end{abstract}


\maketitle
%\tableofcontents


\section{Introduction}

The objective of this paper is to understand the closure of the image of divisor functions. For a complex number $r$, the \textit{divisor function} $\sigma_r : \N \rightarrow \C$ is defined by $$\sigma_r(n) = \sum_{d \mid n} d^{-r},$$ where $\N = \Z_{> 0}.$ 
Since the mid-1980s, various authors have investigated the image $\sigma_r(\N).$ 
In 1986, Laatsch showed that $\sigma_1(\N)$ is dense in $[1, \infty)$ \cite{L86}. In 2000, Weiner asked what can be said about $\overline{\sigma_r(\N)},$ which is the closure of the image of $\sigma_r$ \cite{W00}. To this, in 2015 and 2016, Defant determined various fundamental topological properties of $\overline{\sigma_r(\N)}.$ In particular, he showed that $\overline{\sigma_r(\N)}$ is connected for $r \in \R$ if and only if $r \in (0, \eta],$ where $\eta$ is the unique real number satisfying $$\frac{2^\eta}{2^\eta - 1} \frac{3^\eta + 1}{3^\eta - 1} = \zeta(\eta)$$ and $\zeta$ is the Riemann zeta function \cite{D15}.
In 2017, Sanna proved that when $r > 1,$ the number $C_r$ is finite, giving an algorithm to compute $\overline{\sigma_r(\N)}$ \cite{S17}, and Defant showed that $C_r$ is finite when $\Re(r) > 1$ \cite{D17}.
In 2018, Zubrilina gave explicit lower and upper bounds on $C_r$ \cite{Z18}. In particular, she gave a lower bound of order $\pi(r),$ where $\pi(r)$ is the number of primes less than or equal to $r$, and an upper bound that is exponential in $r$. Additionally, she showed that $C_r$ does not assume all positive integral values. 


In this paper, we introduce a generalized version of the divisor function, defined as follows, and study its properties.

\begin{definition}
    Let $r \in \C,$ and let $K$ be a finite extension of $\Q$. Let $R := \mathcal{O}_K$ be the ring of integers of $K$, and let $I_K$ be the set of integral ideals in $\mathcal{O}_K$. Let $\chi$ be a Dirichlet character.
    Then define the \textit{ideal divisor function on $K$ twisted by $\chi$} as the function $\sigma_{r,K,\chi} : I_K \rightarrow \C$ by $$\sigma_{r,K,\chi}(I) = \sum_{J \mid I} N(J)^{-r}\chi(N(J)),$$ where $N$ denotes the ideal norm. 
% (That is, $N(J) := |R/J|$ is a multiplicative function, and if $J$ is a prime ideal lying over $p\Z$ in $\Z$ and $p$ has inertial degree $f$, then $N(J) = p^j.$) 
Define $C_{r,K,\chi}$ to be the number of connected components of $\overline{\sigma_{r,K,\chi}(I_K)}.$
\end{definition}


We now establish the notation we use throughout the paper. We fix a real number $r > 1$ and a finite  extension $K$ of $\Q.$ For $m \in \N$ and $x \in \R_{>0}$, let $\pi(x)$ denote the number of primes less than or equal to $n$ and $\pi_m(x)$ denote the number of primes less than or equal to $n$ that do not divide $m$. Denote the trivial character by $\chi_0$; note that the function $\sigma_{r, \Q, \chi_0}$ is precisely $\sigma_r.$ In Sections \ref{section-lower-2} and \ref{section-last}, we use $C_{r,K}$ to denote $C_{r,K,\chi_0}.$ In addition, we say a character $\chi$ is \textit{real} if it assumes values in $\{-1,0,1\}.$



This paper is structured as follows. 
In Section \ref{section-finiteness}, we prove that $C_{r,K,\chi}$ is finite, extending Sanna's proof that $C_{r}$ is finite in \cite[Theorem 1.1]{S17}.

\begin{theorem} \label{finiteness}
    Let $r > 1$ be a real number. let $K$ be a finite extension of $\Q,$ and let $\chi$ be a real character. Then $C_{r,K,\chi}$ is finite. 
\end{theorem} 

In Section \ref{section-range}, we prove that when $r$ is fixed and as $\chi$ varies, $C_{r,\Q,\chi}$ assumes all sufficiently large integer values; that is, 

\begin{theorem} \label{rangeC1}
    Let $r > 1$ be a real number. There exists $N \in \N$ such that $ \Z_{\geq N} \subseteq \{C_{r,\Q,\chi} : \chi \text{ is a character}\}.$
\end{theorem}

In particular, we show that when $r$ is large enough, $C_{r,\Q,\chi}$ assumes all positive integral values. 

\begin{theorem} \label{rangeC2}
    There exists a positive real number $r_0$ such that for any $r > r_0,$ we have $\{C_{r,\Q,\chi} : \chi \text{ is a character}\} = \N.$
\end{theorem}

In Section \ref{section-lower-1}, we give a lower bound on $C_{r,\Q,\chi}$ that improves upon the lower bound given by Achenjang and Berger in \cite[Theorem 5.2]{AB19}. 

\begin{theorem} \label{lowerbound1}
    Let $r \geq 3,$ and let $\chi$ be a character of modulus $m$. Then $C_{r,\Q,\chi} \geq 2^{\pi_m(2r-2)-\delta_m},$ where $\delta_m$ equals 0 if $m$ is even and equals 1 otherwise. 
\end{theorem}

Generalizing this result, in Section \ref{section-lower-2}, we give a lower bound on $C_{r,K}$ for finite Galois extensions $K$ of $\Q$.

\begin{theorem} \label{lowerbound2}
    Let $K/\Q$ be a finite Galois extension. For every $\varepsilon > 0,$ there exist constants $\mu_{K,\varepsilon}, r_{K,\varepsilon} > 0$ such that if $r > r_{K, \varepsilon},$ then $$\log C_{r,K} > \mu_{K, \varepsilon}\left( \frac{r^\frac{1}{1+\varepsilon}}{\log r} \right).$$
\end{theorem}

This result implies that fixing $K$, we have $\sup_{r > 1} C_{r,K} = \infty.$ As a complement to this result, in Section \ref{section-last}, we instead fix $s = [K:\Q]$ and $r$ and understand how $C_{r,K}$ behaves as $K$ changes.

\begin{theorem} \label{fix-s}
    Let $r > 1$ be a real number and $s \geq 2$ be an integer. Then $$\sup_{\text{degree-}s \text{ extensions }K/\Q} C_{r,K} = \infty.$$
\end{theorem}

We briefly illustrate some rough intuition for our results.
Below are the plots of the image of the divisor function $\sigma_{2, \Q, \chi}$ on the set $\{1,2,\dots,10^6\},$ where $\chi$ is the principal character $\chi_1$ mod 1 (the constant character) and $\chi_5$ mod 5, respectively.

\begin{figure}[H]
\centering
\includegraphics[scale=0.6]{duluth_graphic_mod1.png}
\caption{The image $\sigma_{2,\Q,\chi_1}(\{1,2,\dots,10^6\}).$}
\includegraphics[scale=0.6]{duluth_graphic_mod5.png}
\caption{The image $\sigma_{2,\Q,\chi_5}(\{1,2,\dots,10^6\}).$}
\end{figure}


As we can reasonably guess from these two figures, the numbers of connected components $C_{2,\Q,\chi}$ are finite. Indeed, they both have 3 connected components; let $I_{1,1}, I_{1,2},$ and $I_{1,3}$ be the connected components for $\chi_1$, and let $I_{5,1},I_{5,2},$ and $I_{5,3}$ be the connected components for $\chi_5.$ We notice that 
\begin{align*}
    I_{i,1} &= \sigma_{2,\Q,\chi_i}(\{n \in \Z_{>0} : 2 \nmid n, 3 \nmid n\})\\
    I_{i,2} &= \sigma_{2,\Q,\chi_i}(\{n \in \Z_{>0} : 2 \nmid n, 3 \mid n\}) \\
    I_{i,3} &= \sigma_{2,\Q,\chi_i}(\{n \in \Z_{>0} : 2 \mid n\})
\end{align*}
for $i = 1, 5.$
In this paper, we will prove this observation that these connected components are determined by the divisibility of the input integer by primes and use it to deduce results about $C_{r,K,\chi}$ in general. 



%%%%%%%%%%%%%%%%%%%%%%%%%%%%%%%%%%%%%%%%%%%%%%%%%%%%%%%%%%%%%%%%%%%%%%%%%%%%%%%%%%%%%%%%%%%%%%%%%%%%%%%%%%%%%%%%%%%%%%%%%%%%%%%%%%%%%%%%%%%%%%%%%%%%%%%%%%%%%%%%%%%%%%%%%%%%%%%%%%%%%%%%%%%%%%%%%%%%%%%%%%%%%%%%%%%%%%%%%%%%%%%%%%%%%%%%%%%%%%%%%%%%%%%%%%%%%%%%%%%%%%%%%%%%%% RESULTS %%%%%%%%%%%%%%%%%%%%
\begin{comment}

\section{Results} 

{\color{MidnightBlue} List of results: 

\begin{itemize}
    \item Finiteness of $C$ \begin{itemize}
        \item known: $C_r$ is finite
        \item If $\chi$ is a real character, then $C_{r,\Q,\chi}$ is finite
        \item (?) If $\chi$ is any character, then $C_{r,\Q,\chi}$ is finite
    \end{itemize}
    \item Range of $C$ \begin{itemize}
        \item For real $r$, range of $C_{r, \chi}$ as $\chi$ varies
    \end{itemize}
    \item Lower bounds on $C$ \begin{itemize}
        \item (written up, though incorrect?) Lower bound on $C_r$ for real $r$ 
        \item Lower bound on $C_{r,\Q,\chi}$ for complex $r$ and character $\chi$ 
    \end{itemize}
\end{itemize}} 


\begin{itemize}
    \item Background results, essentially just adapting everything known for $\Re(r) \geq -1$ in non-character case to character case. Probably can just state? 
    \item $\Re(r) < -1.$ Fix $r$. Then values of $C_{r,\Q,\chi}$ hit like basically? everything as $\chi$ varies.
    \item Finiteness of $C_{r,\Q,\chi}$ for a real character $\chi$.
    \item New lower bound for $C_r$ and thus $C_{r,\Q,\chi}q$

    \item {\color{MidnightBlue} TODO: Write proof of Theorem 7.3. Meanwhile, fill in old lemmata. Then deal with Cebatorev's case.}
\end{itemize}

\end{comment}
\begin{comment}
\section{Notation}

\textcolor{red}{Some notes: 
\begin{itemize}
    \item This is a temporary section, as I didn't want to write my introduction before figuring out the structure. In the same vein, I resorted to clarifying the notation for each section at the beginning of each section for now. But I am also not sure if my notation/the way I present it is the most concise way to do it. 
    \item The biggest change is the inclusion of results with the ideal divisor function, most of which cover the original divisor function. So for example, Section \ref{section-finiteness} proves finiteness for all variants of the divisor function (including number field not $\Q$ and character that's not trivial). But the last two sections, which give lower bounds on $C_{r,\Q,\chi}$ and $C_{r,K,\chi}$ did not get combined---should I?
    \item Section \ref{section-range} seems a little atrocious with regards to notation and structuring, I fear. 
\end{itemize}}

Let $\N$ denote the set of nonnegative integers.


\begin{definition}
    Let $r \in \C$. Then define the \textit{divisor function} $\sigma_r : \N \rightarrow \C$ by $$\sigma_r(n) = \sum_{d \mid n} d^{-r}.$$ Define $C_r$ to be the number of connected components of $\overline{\sigma_r(\N)}.$
\end{definition}

Analogously, we define the following. 

\begin{definition}
    Let $r \in \C$ and $\chi$ be a Dirichlet character modulo $m$. The \textit{divisor function twisted by $\chi$} is defined to be $\sigma_{r,\Q,\chi} : \N \rightarrow \C$ by $$\sigma_r(n) = \sum_{d \mid n} d^{-r} \chi(d).$$ Define $C_{r,\Q,\chi}$ to be the number of connected components of $\overline{\sigma_{r,\Q,\chi}(\N)}.$
\end{definition}

And the following. 

\begin{definition}
    Let $r \in \C,$ and let $K$ be a degree-$s$ abelian extension of $\Q$. Let $R := \mathcal{O}_K$ be the ring of integers of $K$, and let $I_K$ be the set of integral ideals in $\mathcal{O}_K$. Let $\chi$ be a Dirichlet character modulus $m$.
    Then define the \textit{ideal divisor function on $K$ twisted by $\chi$} as the function $\sigma_{r,K,\chi} : I_K \rightarrow \C$ by $$\sigma_{r,K,\chi}(I) = \sum_{J \mid I} N(J)^{-r}\chi(N(J)),$$ where $N$ denotes the ideal norm. That is, $N(J) := |R/J|$ is a multiplicative function, and if $J$ is a prime ideal lying over $p\Z$ in $\Z$ and $p$ has inertial degree $f$, then $N(J) = p^j.$ Define $C_{r,K,\chi}$ to be the number of connected components of $\overline{\sigma_{r,K,\chi}(I_K)}.$
\end{definition}

Note that when $\chi$ is the trivial character, the function $\sigma_{r, \chi, \Q}$ is precisely $\sigma_r.$
\end{comment}
\section{Finiteness of $C_{r,K,\chi}$ for $r > 1$ and real $\chi$} \label{section-finiteness}

The goal of this section is to prove Theorem \ref{finiteness}.
%In this section, we prove the following theorem, which extends \cite{S17, Theorem 1.1}.

\begin{comment}
\begin{theorem} \label{finiteness}
    Let $r > 1,$ let $K$ be a finite abelian extension of $\Q,$ and let $\chi$ be a character that assumes values in $\{-1,0,1\}.$ Then $C_{r,K,\chi}$ is finite. 
\end{theorem}
\end{comment}

Throughout this section, let $\chi$ be a real character with modulus $m$. 
%Denote $\sigma_{r,K,\chi}$ by $\sigma$. 
Order the prime ideals in $\mathcal{O}_K$ by increasing ideal norm; arbitrarily order prime ideals of the same norm. Denote the $i$-th prime ideal in this ordering by $\mathfrak{p}_i.$ 
Let $S_m$ be the finite set $\{i \in \N : \gcd(N(\p_i), m) \neq 1\}$ of indices of prime ideals whose norm are not coprime to $m$. Let $\mathcal{N}_j := \{ I \in I_K : \text{for all }1 \leq i \leq j, \p_i \nmid I\}$ denote the set of ideals whose prime ideal factors have index larger than $j$.
%\textcolor{red}{I introduce similar notation at the beginning of Section \ref{section-range}, just a non-idelic version. I am not sure if I should try to uniformize the notation and put it in the introduction, as using the mathfrak font for integers feels weird.}  

To prove the theorem, we prove a series of useful lemmas. 



%For the remainder of this section, denote $\zeta_{K,i}$ by $\zeta_{K,i}$ and $\zeta_{i,K,m}$ by $\zeta_{K,i}.$ \textcolor{red}


The following lemma establishes a density of norms of certain prime ideals in $\mathcal{O}_K$. Note that it holds for any choice of character $\chi$ with modulus $m$.

\begin{lemma} \label{adj-ideal}
    Consider the prime ideals $\p$ in $\mathcal{O}_K$ such that $\chi(N(\p)) = 1$, and let $\{\p_i'\}_{i \in \N}$ be the sequence of these ideals once ordered by increasing ideal norm. Then $$\lim_{i \rightarrow \infty} \frac{N(\p_{i+1}')}{N(\p_i')} = 1.$$
\end{lemma}

\begin{proof}
    Because there are finitely many prime ideals in $K$ that ramify in $K(\zeta_m),$ it is sufficient to prove the lemma for the sequence $\{\mathfrak p_i'\}_{i \in \N}$ of prime ideals $\p$ in $K$ that are unramified in $K(\zeta_m)$, satisfy $\chi(N(\p)) = 1,$ and are ordered by non-decreasing ideal norm $N$.
    
    Let $F = \Q(\zeta_m)$ and $L = K(\zeta_m).$ 
    Consider the isomorphism
    \begin{align*}
        \text{res} : \Gal(K(\zeta_m)/K) &\ra \Gal(\Q(\zeta_m)/\Q(\zeta_m) \cap K).
    \end{align*} Viewing $\Gal(\Q(\zeta_m)/\Q(\zeta_m) \cap K)$ as a subgroup of $\Gal(\Q(\zeta_m) / \Q) \cong (\Z/ m\Z)^\times,$ the Dirichlet character $\chi$ restricts to a character $$\Tilde{\chi} := \chi \circ \text{res}$$ of $\Gal(K(\zeta_m) / K).$
    For any prime ideal $\mathfrak{p}$ in $K$ that is unramified in $K(\zeta_m)/K$ and lies over $p\Z,$ the map $\res$ sends $\Frob_{\mathfrak{p}, K(\zeta_m)/K}$ to $$\Frob_{p, \Q(\zeta_m)/\Q}^{f_p} \in \Gal(\Q(\zeta_m) / \Q(\zeta_m) \cap K),$$ where $f_p$ is the inertia degree of $\p$ over $p\Z.$ %Then define the character $$\Tilde{\chi} :=\text{res} \circ \chi.$$  
    
    Now, let $H:= \ker \Tilde{\chi}$, which is a nonempty subgroup of $\Gal(K(\zeta_m)/K).$ For any $\sigma \in H,$ consider the prime ideals $\p$ in $K$ that are unramified in $K(\zeta_m)$ and satisfy $\Frob_{\p , K(\zeta_m)/K} = \sigma.$ Ordering these prime ideals by non-decreasing ideal norm $N$, we denote the resulting sequence by  $\{\p_{\sigma, i}\}_{i \in \N}.$ Because $\Tilde{\chi}(\Frob_{\p, K(\zeta_m)/K}) = \chi(N(\p)),$ we see that
    $$\{\p_i'\}_{ i \in \N} = \bigsqcup_{\sigma \in H} \{\p_{\sigma, i}\}_{i \in \N}.$$
    By Chebotarev's density theorem, we have $$\#\{\mathfrak p_{\sigma, i} : N(\mathfrak p_{\sigma,i}) \leq X\} \sim c_\sigma \frac{X}{\log X}$$ for some constant $c_\sigma > 0,$ so $$\lim_{i \ra \infty} \frac{N(\p_{\sigma, i+1})}{N(\p_{\sigma, i})} = 1.$$ As this holds for each $\sigma$ in the finite set $H$, it follows that $$\lim_{i \ra \infty} \frac{N(\p_{i+1}')}{N(\p_{i}')} = 1.$$
    \begin{comment}
    Let $c$ be the conductor of $K$; that is, the minimal integer such that $K \subseteq \Q(\zeta_c).$  Then there exists a nonempty subset $M$ of the congruence classes in $\Z/(mc)\Z$ such that $\chi(x) = 1$ if and only if $x$ lies in a congruence class in $M$. For each rational prime $p$, let $f_p$ be the inertial degree of $p\Z$ in $K$. Recall by the class field theory of $\Q,$ there is a surjection from $(\Z/ c\Z)^*$ to $\text{Gal}(K/\Q)$ that maps the congruence class $a$ (mod $c$) to the Artin symbol of $a\Z.$ Since the primes lying in the same congruence class modulo $mc$ also lie in the same congruence class modulo $c$, they have the same Frobenius element and thus the same inertial degrees. 
    Consider the set of rational primes $p$ whose $f_p$-th power is in a congruence class in $M$; denote the set of congruence classes of these primes modulo $mc$ by $M'.$
    Let $\p$ be a prime ideal in $\mathcal{O}_K$, and let $p\Z$ be the rational prime $\p$ lies over. 
    Then $\chi(N(\p))=1$ if and only if $N(\p) \in M,$ which occurs if and only if $p \in M'.$ 

    For each $x \in M',$ consider the prime ideals $\p$ in $\mathcal{O}_K$ such that $\p$ lies over an rational prime in $x$. Ordering these by increasing ideal norm, denote the resulting sequence of ideals by $\{\p_{x,i}\}_{i \in \N}.$  Since the prime ideals in the sequence $\{\p_{x,i}\}_{i \in \N}$ lie over primes in the same congruence class modulo $c$, they have the same Frobenius element and thus the same inertial degrees. Therefore, by Dirichlet's theorem, for each $x \in M',$ we have $$\lim_{i \rightarrow \infty} \frac{N(\p_{x,i+1})}{N(\p_{x,i})} = 1.$$ Then consider the prime ideals $\p$ in $\mathcal{O}_K$ such that $\p$ lies over an rational prime in an congruence class in $M'.$ Ordering these by increasing ideal norm, the resulting sequence of ideals is precisely $\{\p_{i}'\}_{i \in \N}.$ Then $\{\p_{i}'\}_{i \in \N}$ is, after re-ordering, the union of some ramified prime ideals in $\mathcal{O}_K$---of which there are finitely many---and the sequences $\{\p_{x,i}\}_{i \in \N}$ for $x \in M'.$ Hence, we conclude that $$\lim_{i \rightarrow \infty} \frac{N(\p_{i+1}')}{N(\p_{i}')} = 1.$$ %(Note that if a non-trivial character was not involved, we could more directly apply Cebatorev's density theorem.)
    \end{comment}
\end{proof}

\begin{definition}
    For $r > 1$, a nonnegative integer $i$, and a finite extension $K/\Q,$ define $$\zeta_{K,i}(r) = \prod_{k > i} \frac{1}{1-N(\p_k)^{-r}}.$$ For a positive integer $m$, further define $$\zeta_{K,i,m} (r) = \prod_{k>i; k\not\in S_m} \frac{1}{1-N(\p_k)^{-r}}.$$
\end{definition} 

Let $j_+ = \min\{j \in \mathbb N \setminus S_m : \chi(N(\mathfrak p_j)) = 1\}.$ 
For each integer $j \geq j_+$, let $\mathfrak{q}_j$ be a prime ideal with the largest norm less than or equal to $N(\mathfrak{p}_j)$ such that $\chi(N(\mathfrak{q}_j)) = 1.$  

\begin{lemma} [{\cite[Lemma 2.2]{S17}}] \label{finiteness-finitely-many-j}
    There exist only finitely many $j \in \N \setminus S_m$ satisfying $j \geq j_+$ such that $$\zeta_{K,j,m}(r) < 1 + N(\mathfrak{q}_j)^{-r}.$$
\end{lemma}

\begin{proof}
    Consider the sequence $(y_j)_{j \geq j_+}$ defined by $$y_j := \frac{\zeta_{K,j,m}(r)}{1 + N(\mathfrak{q}_j)^{-r}}.$$ For any $j \geq j_+,$ if $j + 1 \in S_m$, then $y_{j+1}=y_j.$ If $j + 1 \not\in S_m,$ then we compute 
    $$ \frac{y_{j+1}}{y_j} = (1 + N(\mathfrak{q}_j)^{-r}) \cdot \frac{1 - N(\p_{j+1})^{-r}}{1 + N(\mathfrak{q}_{j+1})^{-r}} < (1 + N(\mathfrak{q}_j)^{-r}) \cdot \frac{1 - N(\p_{j+1})^{-r}}{1 + N(\mathfrak{p}_{j+1})^{-r}},$$ because $N(\q_{j+1}) \leq N(\p_{j+1}).$ Now, let $\mathfrak{Q}_j$ be a prime ideal $\mathfrak{Q}$ with the least norm greater than $N(\q_j)$ such that $\chi(N(\QQ)) = 1.$ Then $N(\QQ) \geq N(\p_{j+1}).$ The function from $\R$ to $\R$ defined by $x \mapsto \frac{1-x}{1+x}$ is decreasing, so 
    $$\frac{y_{j+1}}{y_j} < (1 + N(\q_j)^{-r}) \cdot \frac{1 - N(\p_{j+1})^{-r}}{1 + N(\p_{j+1})^{-r}} < (1 + N(\q_j)^{-r}) \cdot \frac{1- N(\QQ_j)^{-r}}{ 1+ N(\QQ_j)^{-r}}.$$
    By Lemma \ref{adj-ideal}, we have $$\lim_{j\rightarrow \infty} \frac{N(\QQ_j)}{N(\q_j)} = 1.$$ 
    Then there exists $j_*$ such that if $j > j_*,$ then $\frac{N(\QQ_j)}{N(\q_j)} < 2^{1/r}$ and thus that $N(\QQ_j)^{-r} > (2 N(\q_j)^r + 1 )^{-1}.$ As a result, for $j > j_*,$ we find 
    $$\frac{y_{j+1}}{y_j} < (1 + N(\q_j)^{-r}) \cdot \frac{1 - (2 N(\q_j)^r + 1 )^{-1}}{1 + (2 N(\q_j)^r + 1 )^{-1}} = 1.$$  Therefore, the sequence $\{y_j\}_{j \geq j_+}$ is eventually non-increasing. Yet $\lim_{j \rightarrow \infty } y_j = 1,$ so $y_j \geq 1$ for all $j > j_*.$ The lemma's statement follows.
\end{proof}
The above lemma allows us to make the following definition.

\begin{definition} \label{defn-j}
    Define $j_{K,\chi}(r)$ to be the least element $j_0$ of $\N \setminus S_m$ such that $j_0 \geq j_+$, 
    \begin{itemize}
        \item[(1)] for all $j \in \N \setminus S_m$ with $j > j_0,$ we have $\zeta_{K,j,m}(r) \geq 1 + N(\mathfrak q_j)^{-r},$ and  
        \item[(2)] for all $j > j_0$ with $\chi(N(\mathfrak p_j)) = -1,$ we have $$ \frac{1}{1-N(\mathfrak p_j)^{-r}} < 1 + N(\mathfrak q_j)^{-r}.$$
    \end{itemize}
\end{definition}
We note that such an element $j_{K,\chi}(r)$ exists. By Lemma \ref{finiteness-finitely-many-j}, sufficiently large $j$ satisfy condition (1). When $\chi(N(\mathfrak p_j)) = -1,$ we have $N(\mathfrak q_j) < N(\mathfrak p_j),$ and the inequality in the second bulleted condition is equivalent to $N(\mathfrak q_j)^r + 1 < N(\mathfrak p_j)^r,$ which holds for sufficiently large $j$.

%For the remander of this section, denote 

For a prime ideal $\p,$ define $$\sigma_{r,K,\chi}(\p^\infty) = \lim_{i \rightarrow \infty} \sigma_{r,K,\chi}(\p^i) = \frac{1}{1 - \chi(N(\mathfrak p)) N(\mathfrak p)^{-r}}.$$

\begin{lemma} [{\cite[Lemma 2.3]{S17}}] \label{finiteness-first-interval}
    For $r>1,$ we have $$\overline{\sigma_{r,K,\chi}( \mathcal{N}_{j_{K,\chi}(r)})} = \left[ \prod_{k > j_{K,\chi}(r); \chi(N(\p_k)) = -1} 1 - N(\p_k)^{-r}, \prod_{k > j_{K,\chi}(r); \chi(N(\p_k)) = 1} \frac{1}{1-N(\p_k)^{-r}}\right].$$
\end{lemma}

\begin{proof}
    Denote the interval above by $[c,d].$
    Take $x \in [c,d].$ We define a sequence $(x_j)_{j \geq j_{K,\chi}(r)}$ as follows. Let $x_{j_{K,\chi}(r)} = c.$ For any $j > j_{K,\chi}(r),$ we define $x_j$ using the following process.
    If $\chi(N(\p_j)) = 0,$ then set $x_j = x_{j-1}.$
    If $\chi(N(\p_j)) = 1,$ then let $a_j$ be the largest nonnegative integer $a$ (including $\infty$) such that $x_{j-1} \cdot \sigma_{r,K,\chi}(\p_j^a) \leq x,$ and let $$x_j = x_{j-1} \cdot \sigma_{r,K,\chi}(\p_j^{a_j}).$$ If $\chi(N(\p_j)) = -1,$ then if $x_{j-1} \cdot \frac{1}{1-N(\p_j)^{-r}} \leq x$, let $b_j = 1$, and let $$x_j = x_{j-1} \cdot \frac{1}{1-N(\p_j)^{-r}}.$$ Otherwise, let $b_j=0$, and let $x_j=x_{j-1}.$ Since the sequence $(x_j)_{j\geq j_{K,\chi}(r)}$ is non-decreasing and bounded, it converges, and let its limit be $\ell.$ Note that $\ell \leq x$. 
    If there are infinitely many $j > j_{K,\chi}(r)$ with $\chi(N(\p_j)) = 1$ and $a_j < \infty$, then $x < x_{j-1} \cdot \frac{1}{1 - N(\p_j)^{-r}}$ for infinitely many $j > j_{K,\chi}(r),$ so $x \leq \ell$ and thus $x = \ell.$
    If there are infinitely many $j >j_{K,\chi}(r)$ with $\chi(N(\p_j)) = -1$ and $b_j=0$, then a similar argument gives $x =\ell.$ Otherwise, there must only be finitely many $j > j_{K,\chi}(r)$ for which $a_j < \infty$ and only finitely many $j > j_{K,\chi}(r)$ for which $b_j = 0.$ 
    Let $j'$ be the least integer greater than equal to $j_{K,\chi}(r)$ such that for each $j > j'$ 
    with $j \not\in S_m,$ the following two conditions hold: when $\chi(N(\p_j)) = 1,$ we have $a_j=\infty,$ and when   $\chi(N(\p_j)) = -1$, we have $b_j = 1.$ Then for any $j > j',$ we have $$x_j = x_{j'} \cdot \prod_{j' < k \leq j; k \not\in S_m} \frac{1}{1-N(\p_k)^{-r}}$$ and thus that $$\ell = x_{j'} \cdot \prod_{k > j'; k \not\in S_m} \frac{1}{1-N(\p_k)^{-r}}.$$ If $j' = j_{K,\chi}(r),$ then $$\ell = c \cdot \prod_{k > j_{K,\chi}(r); k \not\in S_m} \frac{1}{1-N(\p_k)^{-r}} = d$$ and $\ell \leq x \leq d,$ so $\ell = x.$
    Otherwise, if $j' > j_{K,\chi}(r),$ then by the minimality of $j'$, we have $j' \not\in S_m.$ Thus, either $\chi(N(\p_{j'})) = 1$ and $a_{j'} < \infty,$ or $\chi(N(\p_{j'})) = -1$ and $b_{j'} = 0.$ 
    If $a_{j'} <\infty,$ then \begin{align*}
        x_{j'} \cdot \prod_{k > j'; k \not\in S_m} \frac{1}{1-N(\p_k)^{-r}} = \ell \leq x &< x_{j'-1} \cdot \sigma_{r,K,\chi}(\p_{j'}^{a_{j'} + 1}) \\
        &\leq \frac{x_{j'}}{\sigma_{r,K,\chi}(\p_{j'}^{a_{j'}})} \cdot  \sigma_{r,K,\chi}(\p_{j'}^{a_{j'} + 1}) \\
        &\leq x_{j'}  (1 + N(\p_{j'})^{-r})\\
        &= x_{j'}(1 + N(\q_{j'})^{-r}),
    \end{align*} where the last equality holds because $\chi(N(\p_{j'})) = 1.$ This contradicts condition (1) in the definition of $j_{K,\chi}(r).$ Similarly, if $b_{j'} = 0,$ then \begin{align*}
        x_{j'} \cdot \prod_{k > j'; k \not\in S_m} \frac{1}{1-N(\p_k)^{-r}} = \ell \leq x &< x_{j'-1} \cdot \frac{1}{1-N(\p_{j'})^{-r}} \\
        &= x_{j'} \cdot \frac{1}{1-N(\p_{j'})^{-r}} \\
        &< x_{j'} (1 + N(\q_{j'})^{-r}),
    \end{align*} where the last inequality holds due to the condition (2) in the definition of $j_{K,\chi}(r).$ This contradicts condition (1) in the definition of $j_{K,\chi}(r) \leq j'.$ Hence, $j' =j_{K,\chi}(r),$ so $x = \ell.$
\end{proof}





\begin{lemma} [{\cite[Lemma 2.4]{S17}}] \label{finiteness-induct}
    For any $1 \leq j \leq j_{K,\chi}(r),$  $$\overline{\sigma_{r,K,\chi}(\mathcal{N}_{j-1})} = \bigcup_{a \in \Z_{\geq 0} \cup \{\infty\}} \sigma_{r,K,\chi}(\p_j^a) \cdot \overline{\sigma_{r,K,\chi}(\mathcal{N}_j)}.$$
\end{lemma}

\begin{proof}
    It is clear that the right-hand side is contained in the left-hand side. To see the reverse containment, take $x \in \overline{\sigma_{r,K,\chi}(\mathcal{N}_{j-1})}.$ Then there exists a sequence $\{I_i\}_{i \in \N} \subset \mathcal{N}_{j-1}$ such that $\lim_{i \rightarrow \infty} \sigma_{r,K,\chi}(I_i) = x.$ For each $i > 0,$ let $a_i$ be the largest power of $\p_j$ dividing $I_i.$ If $\{a_i\}_{i \in \N}$ is bounded, then there exists a nonnegative integer $a_0$ that is attained infinitely many times in this sequence; let $\{i_\ell\}_{\ell \in \N}$ be the corresponding subindices. Then $$\frac{x}{\sigma_{r,K,\chi}(\p_j^{a_0})} = \lim_{\ell \rightarrow \infty} \sigma_{r,K,\chi}\left( I_{i_\ell} \p_j^{-a_0}\right) \in \overline{\sigma_{r,K,\chi}(\mathcal{N}_j)}.$$ If $\{a_i\}_{i \in \N}$ is unbounded, there exists a subsequence $\{a_{i_\ell}\}_{\ell \in \N}$ that converges to $\infty.$ Then $$\frac{x}{\sigma_{r,K,\chi}(\p_j^\infty)} = \lim_{\ell \rightarrow \infty} \sigma_{r,K,\chi}(I_{i_\ell} \p_j^{-a_{i_\ell}}) \in \overline{\sigma_{r,K,\chi}(\mathcal{N}_j)}.$$
\end{proof}

\begin{lemma} [{\cite[Lemma 2.5]{S17}}] \label{finiteness-last}
    Let $j > 0,$ and let $0 < \alpha < \beta.$  The set $$\bigcup_{a \in \Z_{\geq 0} \cup \{\infty\}} \sigma_{r,K,\chi}(\p_j^a) \cdot [\alpha,\beta]$$ is a finite 
    union of closed bounded intervals. In particular, if $\chi(N(\mathfrak p_j)) = 1$ and  $a_0$ is the least nonnegative integer such that $$\frac{\sigma_{r,K,\chi}(\p_j^{a_0+1})}{\sigma_{r,K,\chi}(\p_j^{a_0})} \leq \frac{\beta}{\alpha},$$ then the set above is the finite union of precisely $a_0+1$ pairwise disjoint closed bounded intervals. 
\end{lemma}

\begin{proof}
    We consider the three possible cases of the value of $\chi(N(\mathfrak p_j)).$ If $\chi(N(\mathfrak p_j)) = 0,$ then the desired union is simply $[\alpha,\beta].$ Next, suppose $\chi(N(\mathfrak p_j)) \neq 0.$ Because $\sigma_{r,K,\chi}(\mathfrak p_j^a)$ converges to $\sigma_{r,K,\chi}(\mathfrak p_j^\infty)$ as $a$ tends to $\infty,$ there exists $A$ for which $a \geq A$ implies the intervals $\sigma_{r,K,\chi}(\mathfrak p_j^a)[\alpha,\beta]$ and $\sigma_{r,K,\chi}(\mathfrak p_j^{a+1})[\alpha,\beta]$ intersect. The total union can thus be express as the union of at most $A+1$ pairwise disjoint closed bounded intervals. In particular, when $\chi(N(\mathfrak p_j)) = 1,$ if we take $A$ to be $a_0,$ then the total union consists of precisely $a_0+1$ pairwise disjoint closed bounded intervals. 
\end{proof}

We are ready to prove Theorem \ref{finiteness}.

\begin{proof}[Proof of Theorem \ref{finiteness}]
    By Lemma \ref{finiteness-first-interval}, $\overline{\sigma_{r,K,\chi}(\mathcal{N}_{j_{K,\chi}(r)})}$ is a single bounded closed interval. We use downward induction to show that for each $1 \leq j \leq j_{K,\chi}(r),$ the closure $\overline{\sigma_{r,K,\chi}(\mathcal{N}_{j-1})}$ is a finite union of pairwise disjoint closed bounded intervals. By Lemmas \ref{finiteness-induct} and \ref{finiteness-last}, the claim holds for $j = j_{K,\chi}(r).$ Suppose the claim holds for some $j = j'$; that is, $\overline{\sigma_{r,K,\chi}(\mathcal{N}_{j'})} = \bigsqcup_{t \in T} C_t$, where $T$ is finite and each $C_t$ is a bounded closed interval. By Lemma \ref{finiteness-induct}, we have $$\overline{\sigma_{r,K,\chi}(\mathcal{N}_{j'-1})} = \bigcup_{t \in T} \left( \bigcup_{a \in \Z_{\geq 0} \cup \{\infty\}} \sigma_{r,K,\chi}(\p_{j'}^a) \cdot C_t \right).$$ For each $t \in T,$ the inner union is a finite union of closed bounded intervals by Lemma \ref{finiteness-last}, so $\overline{\sigma_{r,K,\chi}(\mathcal{N}_{j'-1})}$ is the finite union of closed bounded intervals. Since $\mathcal{N}_0 = I_K,$ the theorem is proved.
\end{proof}





\begin{comment}
\section{lemmas for Principal !!}

\textcolor{violet}{Remember that these lemmas work only for principal characters; $\chi_m$ is the principal character modulo $m$!}

\begin{definition}
    \textcolor{blue}{(TODO: will move to Notation)} For $r > 1$ and a nonnegative integer $i$, define $$\zeta_i(r) = \prod_{k > i} \frac{1}{1-p_k^{-r}}.$$ For a positive integer $m$, further define $$\zeta_{i,m} (r) = \prod_{k>i; k\not\in S_m} \frac{1}{1-p_k^{-r}}.$$
\end{definition}

\begin{lemma}[\cite{S17}, Lemma 2.2]
    For $r > 1,$ there exist only finitely many $j \in \N\setminus S$ such that $$\zeta_{j, m}(r) < 1 + p_j^{-r}.$$ \textcolor{red}{follows from Lem 5.1}
\end{lemma}


\begin{definition}
    Define $j_m$ to be the least element of $\N \setminus S$ such that for all $j \in \N \setminus S$ with $j>j_m,$ we have $$\zeta_{K,j,m} \geq 1 + p_j^{-r}.$$
\end{definition}

\begin{lemma}[\cite{S17}, Lemma 2.3] \label{sanna2.3}
    For $r > 1,$ we have $$\overline{\sigma_{r, \chi_m}(\mathcal{N}_{j_m})} = \left[ 1, \zeta_{j_m, m}(r)\right].$$ \textcolor{red}{follows from Lem 5.3}
\end{lemma}

\begin{lemma}[\cite{S17}, Lemma 2.4] \label{sanna2.4}
    For each positive integer $j \leq j_m$, we have $$\overline{\sigma_{-r,\chi_m}(\mathcal{N}_{j-1})} = \bigcup_{a \in \N \cup \{\infty\}} \sigma_{-r,\chi_m}(p_j^a) \cdot \overline{\sigma_{-r,\chi_m}(\mathcal{N}_{j})}.$$ 
\end{lemma}





\begin{lemma}[\cite{S17}, Lemma 2.5] \label{sanna2.5}
    Given $\beta>\alpha\geq1$ and $j \in \N^+\setminus S,$ let $a_0$ be the least nonnegative integer such that $$\frac{\sigma_{-r,\chi_m}(p_j^{a_0+1})}{\sigma_{-r,\chi_m}(p_j^{a_0})} \leq \frac{\beta}{\alpha}.$$ For each nonnegative integer $a < a_0,$ define $J_a := \sigma_{-r,\chi_m}(p_j^{a}) \cdot [\alpha,\beta],$ and define $J_{a_0} := [\sigma_{-r,\chi_m}(p_j^{a_0})\alpha, \sigma_{-r,\chi_m}(p_j^{\infty})\beta].$ Then $J_0,\ldots,J_{a_0}$ are pairwisely disjoint, and we have $$\bigcup_{a \in \N \cup \{\infty\}} \sigma_{-r,\chi_m}(p_j^a) \cdot [\alpha,\beta] = J_0 \cup \cdots \cup J_{a_0}.$$
\end{lemma}

\begin{theorem}
    If $\chi$ assumes values in $\{0,1\},$ then $C_{r,\Q,\chi}$ is finite. 
\end{theorem}

\textcolor{green}{I wonder if you can just condense this into the next section?}


\section{for real character}

Let $\chi$ be a character with modulus $m$ that takes values in $\{-1,0,1\}.$ For each positive integer $j$, let $q_j$ be the largest prime $p$ less than or equal to $p_j$ such that $\chi(p) = 1.$

\begin{lemma} \label{lemma1-finite}
    For $r>1,$ there exist only finitely many $j \in \N \setminus S_m$ such that $$\zeta_{K,j,m}(r) < 1 + q_j^{-r}.$$
\end{lemma}

\begin{proof}
    Consider the sequence $(y_j)_{j \in \N \setminus S_m}$ defined by $$y_j := \frac{\zeta_{j, m}(r)}{1+q_j^{-r}}.$$
    For any $j \in \N \setminus S_m,$ we compute $$\frac{y_{j+1}}{y_j} = (1 + q_j^{-r}) \cdot \frac{1 - p_{j+1}^{-r}}{1 + q_{j+1}^{-r}} < (1 + q_j^{-r}) \cdot \frac{1 - p_{j+1}^{-r}}{1 + p_{j+1}^{-r}}$$ because $q_{j+1} \leq p_{j+1}.$ Now, let $Q_j$ be the least prime greater than $q_j$ such that $\chi(Q) = -1$. Then $Q \geq p_{j+1}.$
    The function from $\R$ to $\R$ defined by $x \mapsto \frac{1-x}{1+x}$ is increasing, so $$\frac{y_{j+1}}{y_j} < (1 + q_j^{-r}) \cdot \frac{1 - p_{j+1}^{-r}}{1 + p_{j+1}^{-r}} < (1 + q_j^{-r}) \cdot \frac{1- Q_j^{-r}}{ 1+ Q_j^{-r}}.$$ Recall that there exists some nonempty subset $M$ of $(\Z/m\Z)^*$ such that $\chi(p) = 1$ if and only if $p \mod{m} \in M.$
    Now, by Dirichlet's theorem on primes in arithmetic progressions, for any $a \in (\Z/m\Z)^*,$ the ratio of the $(k+1)$-th to the $k$-th prime that is equivalent to $a$ modulo $m$ approaches 1 as $k$ approaches $\infty,$ so $\lim_{j \rightarrow \infty} \frac{Q_j}{q_j} = 1$ \textcolor{red}{(isn't this just equality?)}. Then there exists $j_0$ such that if $j > j_0,$ then $\frac{Q_j}{q_j} < 2^{1/r}$ and thus that $Q_j^{-r} > (2q_j^{r} + 1)^{-1}.$ As a result, for $j > j_0,$ we have $$\frac{y_{j+1}}{y_j} < (1 + q_j^{-r}) \cdot \frac{1- (2q_j^r+1)^{-1}}{ 1+ (2q_j^r+1)^{-1}} = 1.$$ Therefore, the sequence $(y_j)$ is eventually strictly decreasing, yet $\lim_{j \rightarrow \infty} y_j = 1,$ so $y_j > 1$ for all $j >j_0.$ The statement follows. 
\end{proof}

\begin{definition}
    Define $j_m$ to be the least element of $\N \setminus S_m$ such that for all $j \in \N \setminus S_m$ with $j > j_m,$ we have $$\zeta_{K,j,m}(r) \geq 1 + p_j^{-r}.$$
\end{definition}

\begin{lemma}
    For $r>1,$ we have $$\overline{\sigma_{r,\Q,\chi}( \mathcal{N}_{j_m})} = \left[ \prod_{k > j_m; \chi(p_k) = -1} 1 - p_k^{-r}, \prod_{k > j_m; \chi(p_k) = 1} \frac{1}{1-p_k^{-r}}\right].$$
\end{lemma}

\begin{proof}
    Denote the interval above by $[c,d].$
    Take $x \in [c,d].$ We define a sequence $(x_j)_{j \geq j_m}$ as follows. Let $x_{j_m} = c.$ For any $j > j_m,$ if $\chi(p_j) = 1,$ then let $a_j$ be the largest nonnegative integer $a$ (including $\infty$) such that $x_{j-1} \cdot \sigma(p_j^a) \leq x,$ and let $x_j = x_{j-1} \cdot \sigma(p_j^{a_j}).$ If $\chi(p_j) = -1,$ then if $x_{j-1} \cdot \frac{1}{1-p_j^{-r}} \leq x$, let $b_j = 1$, and let $x_j = x_{j-1} \cdot \frac{1}{1-p_j^{-r}}.$ Otherwise, let $b_j=0$, and let $x_j=x_{j-1}.$ Since the sequence $(x_j)_{j\geq j_m}$ is non-increasing and bounded, it converges, and let its limit be $\ell.$ Note that $\ell \leq x$. If $a_j <\infty$ for infinitely many $j > j_m$, then $x < x_{j-1} \cdot \frac{1}{1 - p_j^{-r}}$ for infinitely many $j > j_m,$ so $x \leq \ell$ and thus $x = \ell.$
    If $b_j =0$ for infinitely many $j >j_m$, then a similar argument gives $x =\ell.$ Otherwise, there must only be finitely many $j > j_m$ for which $a_j < \infty$ and only finitely many $j > j_m$ for which $b_j = 0.$ Let $j'$ be the least integer greater or equal to $j_m$ such that $a_j = \infty$ and $b_j=1$ for all $j >j'.$ Then for any $j > j',$ we have $$x_j = x_{j'} \cdot \prod_{k > j'; k \not\in S} \frac{1}{1-p_k^{-r}}$$ and thus that $$\ell = x_{j'} \cdot \prod_{k > j'; k \not\in S} \frac{1}{1-p_k^{-r}}.$$ If $j' = j_m,$ then $$\ell = c \cdot \prod_{k > j_m; k \not\in S} \frac{1}{1-p_k^{-r}} = d$$ and $\ell \leq x \leq d,$ so $\ell = x.$
    Otherwise, we have either $a_{j'} < \infty$ or $b_{j'} = 0.$ If $a_{j'} <\infty,$ then \begin{align*}
        x_{j'} \cdot \prod_{k > j'; k \not\in S} \frac{1}{1-p_k^{-r}} = \ell \leq x &\leq x_{j'-1} \cdot \sigma_r(p_{j'}^{a_{j'} + 1}) \\
        &\leq \frac{x_{j'}}{\sigma_r(p_{j'}^{a_{j'}})} \cdot  \sigma_r(p_{j'}^{a_{j'} + 1}) \\
        &\leq x_{j'}  (1 + p_{j'}^{-r}),
    \end{align*} which contradicts the definition of $j_m \leq j'.$ Similarly, if $b_{j'} = 0,$ then \begin{align*}
        x_{j'} \cdot \prod_{k > j'; k \not\in S} \frac{1}{1-p_k^{-r}} = \ell \leq x &\leq x_{j'-1} \cdot \frac{1}{1-p_{j'}^{-r}} \\
        &= x_{j'} \cdot \frac{1}{1-p_{j'}^{-r}} \\
        &< x_{j'} (1 + q_{j'}^{-r})
    \end{align*} because $q_{j'} +1 \leq p_{j'}$. This again contradicts the definition of $j_m \leq j'.$ Hence, $j' =j_m,$ so $x = \ell.$
\end{proof}

\textcolor{blue}{more lemmata}

\begin{theorem}
    $C_{r,\Q,\chi}$ is finite for $r > 1.$ 
\end{theorem}
\end{comment}

\section{Range of $C_{r,\Q,\chi}$ when $r$ is fixed} \label{section-range}

%\textcolor{red}{[Probably can modify this to work for any $K$ as well.]}

\begin{comment}
\textcolor{blue}{If $1 < r < 1.94$-ish, then it is fairly obvious that $C_{r,\Q,\chi}$ tends to infinity, where $\chi$'s modulus keeps getting multiplied by then next prime and all the primes up to some $q$ are mod-ed out, where $\chi(q) = 1.$ \textbf{Define} what it means for modulus to ``keep getting multiplied." \textbf{To-do list:}
\begin{itemize}
    \item I have a proof for $j_0<2$ case, but need to check if $C = 2$ situation still works (does not work) and figure out how to adjust notation, since the last segment on $C - C \leq 1$ is the same argument. 
    So we can do: (most of this requires $r$ to be real) notation prior to all of this. A proposition for last, shared segment. A separate theorem for saying when $j_0<2$ (certain $r$ values), $C$ attains all but finitely many (uses that first prop). [A proposition on how ``multiplying the modulus" of any, potentially complex characters increases $C.$ Here, $r$ can be complex.] Then conclude with a theorem for when $j_0 \geq 2.$
    \item Fix Lemma 6.1; proof is not right?
    \item Rewrite Theorem 6.2. Separate by closure instead; be cautious of how you're defining $A$, square root number of primes?
    \item Split up proof of main theorem?
    \item Write up proofs for Sanna-attributed Lemmas. 
    \item Write proofs for finiteness with the real character.
\end{itemize}}
\end{comment}


In this section, we prove some results about the values of $C_{r,\Q,\chi}$ for $r > 1$ and a Dirichlet character $\chi.$ In \cite[Theorem 3.3]{D16}, Defant shows that for $r \in \C$ with $\Re(r) > 1,$ we have $\lim_{r \rightarrow \infty} C_{r,\Q,\chi_0} = \infty.$ One can deduce from the theorem's proof that for a fixed character $\chi,$ we have $\lim_{r \rightarrow \infty} C_{r,\Q,\chi} = \infty.$ In this vein, we fix $r>1$ and determine what values $C_{r,\Q,\chi}$ can attain for arbitrary real character $\chi.$ 

In particular, we will prove Theorem \ref{rangeC1}, which asserts that for any $r > 1,$ there exists large enough $N \in \mathbb{N}$ for which $C_{r,\Q,\chi}$ can attain any integral value at least $N$. 
The key idea of the proof is that increasing the number and the magnitude of the distinct primes dividing a character's modulus forces additional jumps in the divisor function's range, thus generally increasing the number of connected components in the range. Thus, the most natural collection of characters we can use to witness and control these effects in the range are principal characters, which are determined entirely by the set of distinct prime factors of their moduli. 


In this section, we use notation in a similar fashion to that of the previous section. Fix $r > 1.$ Let $p_j$ be the $j$-th rational prime and $\mathcal{N}_j$ be the set of positive integers not divisible by the first $j$ primes. For any positive integer $m$, let $\chi_{m}$ be the principal modulus-$m$ character. Lastly, recalling Definition \ref{defn-j}, set $j_m:= j_{\Q,\chi_m}(r).$

We proceed to define some new notation to select the characters $\chi$ in order for $C_{r,\Q,\chi}$ to attain any sufficiently large integral value. We define $\ell := \max(2, j_1)$ and 
$$ i_0= \begin{cases}
    j_1 & \text{if }j_1 \geq 2 \\
    \min\bigl\{i\in \mathbb{Z}_{>0} : \zeta_{\Q, i}(r) < 1+3^{-r}\} & \text{else.}
\end{cases}$$
Note that $i_0 \geq \ell,$ so we finally define the sequence of positive integers $(m_i)_{i \geq i_0}$ by $$m_i := \frac{1}{p_\ell} \prod_{p \leq p_i} p.$$ The moduli $m_i$ are chosen so that $p_\ell$ is the unique prime that contributes Euler factors to the image of $\sigma_{r,\Q,\chi}$. In particular, the prime $p_\ell$ creates ``copies" of $[1, \zeta_{\Q,i}(r)]$ in $\overline{\sigma_{r,\Q,\chi}(\N)}$. Increasing $i$ shrinks the lengths of these interval copies and thus increases the number of connected components.

\begin{lemma} \label{rangeClemma}
    Let $i \geq i_0.$ Then $(C_{r,\Q,\chi_{m_i}})_{i \geq i_0}$ is a non-decreasing sequence that is not eventually constant. In particular, for any $i \geq i_0,$ we have $$C_{r,\Q,\chi_{m_i}} = \left \lceil \log_{p_\ell^{-r}} \frac{\zeta_{\Q,i}(r)-1}{\zeta_{\Q,i}(r) - p_\ell^{-r}} \right \rceil.$$
\end{lemma}


\begin{proof}
    Let $i \geq i_0.$ We first show that $j_{m_i} = \ell.$ It is clear that $j_+ = \ell.$ Now, take any $j \in \mathbb N \setminus S_{m_i}$ for which $j > \ell.$ Because $m_i$ is divisible by each $p_k$ with $k \leq i$ except $p_\ell,$ this implies $j > i.$ Hence, $\zeta_{\Q, j , m_i}(r) = \zeta_{\Q,j}(r).$ For any $j  > i,$ we have $\zeta_{\Q,j}(r) \geq 1 + p_j^{-r}$ because $i \geq i_0.$ Thus, $\ell$ satisfies the defining property of $j_{m_i}.$ 

    Since $\sigma_{r,\Q,\chi_{m_i}} (p_j^a) = 1$ for $1 \leq j \leq i$ with $j \neq \ell,$ we apply Lemma \ref{finiteness-induct} for $1 \leq j \leq \ell$ and subsequently apply Lemma \ref{finiteness-first-interval}
    to obtain \begin{align*}
        \overline{\sigma_{r, \Q,\chi_{m_i}} (\N)} &= \overline{\sigma_{r, \Q,\chi_{m_i}} (\mathcal{N}_1)} \\
        &= \bigcup_{a \in \Z_{\geq 0} \cup\{\infty\}} \sigma_{r, \Q,\chi_{m_i}} (p_\ell^a) \cdot \overline{\sigma_{r, \Q,\chi_{m_i}} (\mathcal{N}_{\ell})} \\
        &= \bigcup_{a \in \Z_{\geq 0} \cup\{\infty\}} \sigma_{r, \Q,\chi_{m_i}} (p_\ell^a) \cdot [1, \zeta_{\Q, j_{m_i}, m_i}(r)].
    \end{align*} Because $j_{m_i} = \ell,$ we have $\zeta_{\Q, j_{m_i}, m_i}(r) = \zeta_{\Q, \ell, m_i}(r).$ As each prime $p_k$ with $\ell < k \leq i$ divides $m_i,$ we have $\zeta_{\Q,\ell,m_i}(r) = \zeta_{\Q,i}(r).$ Thus, we have $$\overline{\sigma_{r, \Q,\chi_{m_i}} (\N)} = \bigcup_{a \in \Z_{\geq 0} \cup \{\infty\}}  \sigma_{r, \Q,\chi_{m_i}} (p_\ell^a) \cdot [1, \zeta_{\Q, i}(r)].$$
    %Before proceeding, we note that $j_{m_i} \leq i-1$. Because $i \geq h > j_1,$ we have $\zeta_{\Q, i , m_i}(r) = \zeta_{\Q, i}(r) \geq 1 + p_i^{-r},$ meaning $j_{m_i} \leq i-1.$ Then applying Lemma \ref{finiteness-induct} to $\overline{\sigma_{r, \Q,\chi_{m_i}} (\mathcal{N}_{\ell})}$ for $\ell+1 \leq j \leq j_{m_i}$ and subsequently applying Lemma \ref{finiteness-first-interval} gives us $$\overline{\sigma_{r, \Q,\chi_{m_i}} (\N)} = \bigcup_{a \in \N \cup\{\infty\}} \sigma_{r, \Q,\chi_{m_i}} (p_\ell^a) \cdot \overline{\sigma_{r, \Q,\chi_{m_i}} (\mathcal{N}_{j_{m_i}})} = \bigcup_{a \in \N \cup\{\infty\}} \sigma_{r, \Q,\chi_{m_i}} (p_\ell^a) \cdot [1, \zeta_{\Q, j_{m_i}, m_i}(r)].$$ 
    %We show that $\zeta_{\Q, j_{m_i},m_i}(r) = \zeta_{\Q,i}(r).$ To do so, we verify that for $i \geq h,$ we have $ \max(2,j_1) \leq j_{m_i} \leq i-1.$ The claim would then follow, as $k \not\in S_{m_i}$ is equivalent to $k =\ell$ or $k > i.$ We have shown earlier that $j_{m_i} \leq i-1.$ We first check $j_{m_i} \geq j_1.$ For any $j \geq j_1,$ we have $\zeta_{\Q, j, m_i}(r) < \zeta_{\Q,j}(r) < 1+p_j^{-r},$ so $j_1 \leq j_{m_i}$. We proceed to check $j_{m_i} \geq 2.$ If $j_1 \geq 2,$ then this is clear. If $j_1 < 2,$ then for $i \geq h,$ we have $\zeta_{\Q,2, m_i}(r) = \zeta_{\Q,i}(r) < \zeta_{\Q,h}(r) < 1 + 3^{-r},$ so $j_{m_i} \geq 2.$ 

    Now, let $a_i$ be the least nonnegative integer such that \begin{equation} \label{eq1}
        \frac{ \sigma_{r,\Q,\chi_{m_i}} (p_{\ell}^{a_i+1}) }{ \sigma_{r,\Q,\chi_{m_i}} (p_{\ell}^{a_i}) } \leq \zeta_{\Q,i}(r).
    \end{equation} By Lemma \ref{finiteness-last}, we have $C_{r,\Q,\chi_{m_i}} = a_i+1.$
    Note that the left hand side of the equation is $$\frac{ 1+p_{\ell}^{-r} + \cdots + p_{\ell}^{-(a_i+1)r} }{  1+p_{\ell}^{-r} + \cdots + p_{\ell}^{-a_i r}  } = \frac{1 - (p_{\ell}^{-r})^{a_i+2}}{1 - (p_{\ell}^{-r})^{a_i+1}}.$$ Algebraic manipulation shows that $$C_{r,\Q,\chi_{m_i}} = a_i +1 = \left\lceil \log_{p_\ell^{-r}} \frac{\zeta_{\Q,i}(r) - 1}{\zeta_{\Q,i}(r) - p_\ell^{-r}} -1 \right\rceil+1 = \left\lceil \log_{p_\ell^{-r}} \frac{\zeta_{\Q,i}(r) - 1}{\zeta_{\Q,i}(r) - p_\ell^{-r}} \right\rceil.$$ As $i$ increases, the expression inside the ceiling increases and approaches $\infty$ because $\zeta_{\Q,i}(r)$ approaches 1 from the right. Thus, $(C_{r,\Q,\chi_{m_i}})_{i \geq i_0}$ is non-decreasing and not eventually constant.
\end{proof}







We proceed to prove Theorem \ref{rangeC1}. 

\begin{proof}[Proof of Theorem \ref{rangeC1}]
    For each $i$, define $y_i = \zeta_{\Q,i}(r).$
    By Lemma \ref{rangeClemma}, it suffices to prove that $C_{r,\Q,\chi_{m_{i+1}}} - C_{r,\Q,\chi_{m_i}} \leq 1$ for any $ i \geq i_0.$  Recall from Lemma \ref{rangeClemma} that $C_{r,\Q,\chi_{m_i}} = a_i+1,$ where $a_i$ is the least nonnegative integer satisfying Equation (\ref{eq1}).
    We essentially show that the sequence $$\left(\frac{1 - (p_\ell^{-r})^{k+2}}{1 - (p_\ell^{-r})^{k+1}} \right)_{k \geq 0}$$ converges to 1 at a faster rate than $(\zeta_{\Q,i}(r))_{i \geq i_0}$ does. 
    In particular, we prove that for any $k \geq0$ and $i \geq i_0,$ we have $$\frac{\frac{1 - (p_\ell^{-r})^{k+3}}{1 - (p_\ell^{-r})^{k+2}}  - 1}{\frac{1 - (p_\ell^{-r})^{k+2}}{1 - (p_\ell^{-r})^{k+1}}  - 1} < \frac{\zeta_{\Q,i+1}(r) - 1}{\zeta_{\Q,i}(r)- 1}.$$ 
    This proves the desired, because if $C_{r,\Q, \chi_{m_{i+1}}} - C_{r,\Q, \chi_{m_i}} > 1$ for some $i \geq i_0,$ then $a_{i+1} \geq a_i+ 2.$ This would imply that there exists $k \geq 0$ such that $$\frac{1 - (p_\ell^{-r})^{k+2}}{1 - (p_\ell^{-r})^{k+1}} \leq y_{i}\text{ and } \frac{1 - (p_\ell^{-r})^{k+3}}{1 - (p_\ell^{-r})^{k+2}} > y_{i+1},$$ which would contradict the claim. 
    
    For any $k \geq 0,$ we find that \begin{align*}
        \frac{ \frac{1 - (p_\ell^{-r})^{k+3}}{1 - (p_\ell^{-r})^{k+2}} - 1}{\frac{1 - (p_\ell^{-r})^{k+2}}{1 - (p_\ell^{-r})^{k+1}} - 1} 
        &= p_\ell^{-r} \cdot \frac{1 - (p_\ell^{-r})^{k+1} }{1 - (p_\ell^{-r})^{k+2}} \\
        &< 3^{-r}
    \end{align*} because $\ell \geq 2.$
    On the other hand, for any $i \geq i_0,$ we find that 
    \begin{align*}
        \frac{\zeta_{\Q,i+1}(r) - \zeta_{\Q,i+2}(r)}{\zeta_{\Q,i}(r) - \zeta_{\Q,i+1}(r)} &= \frac{ \left( \frac{1}{1-p_{i+2}^{-r}} -  1 \right) \zeta_{\Q,i+2}(r) }{  \left( \frac{1}{1-p_{i+1}^{-r}} -  1 \right) \zeta_{\Q,i+1}(r)  }\\
        &= \frac{p_{i+2}^{-r} (1 - p_{i+1}^{-r})}{p_{i+1}^{-r}} \\
        &> 2^{-r}(1 - 3^{-r}),
    \end{align*}
    where the last step follows from Bertrand's postulate and the fact that $p_{i+1} \geq p_2 = 3$. Inductively, we conclude that for any $k \geq 0,$ we have $$\zeta_{\Q,i+k}(r) - \zeta_{\Q,i+k+1}(r) > (2^{-r} (1 - 3^{-r}) )^k (\zeta_{\Q,i}(r) - \zeta_{\Q,i+1}(r)) > (3^{-r} )^k (\zeta_{\Q,i}(r) - \zeta_{\Q,i+1}(r)).$$ As a result, \begin{align*}
        \zeta_{\Q,i+1}(r) - 1 &= \sum_{k=1}^\infty \zeta_{\Q,i+k}(r) - \zeta_{\Q,i+k+1}(r) \\&> (3^{-r} + (3^{-r})^2 + \cdots ) (\zeta_{\Q,i}(r) - \zeta_{\Q,i+1}(r)) \\ &= \frac{1}{3^r - 1} (\zeta_{\Q,i}(r) - \zeta_{\Q,i+1}(r)),
    \end{align*}so we find that $$\frac{\zeta_{\Q,i+1}(r) - 1}{\zeta_{\Q,i}(r)- 1} > \frac{ \frac{1}{3^r - 1} }{ \frac{1}{3^r - 1} + 1} = 3^{-r}.$$ 
    Then setting $N$ in the theorem statement to be $C_{r,\Q,\chi_{m_{i_0}}}$ gives the theorem.
\end{proof}

We proceed to prove Theorem \ref{rangeC2}, where we show that when $r$ is large enough, we can describe $C_{r,\Q,\chi}$ in greater detail.





\begin{proof}[Proof of Theorem \ref{rangeC2}]
    We first show the existence of a real number $r_0 > 1$ such that $$\sum_{n \geq 3} n^{-r} < 2^{-r} \text{ and } \sum_{n\geq 4} n^{-r} < 3^{-r}$$ hold for any $r > r_0.$ Indeed, we can bound $$\sum_{n \geq 3} n^{-r} \leq 3^{-r} + \int_3^\infty x^{-r} \; dx = 3^{-r} \left(1 + \frac{3}{r-1} \right),$$ which is less than $2^{-r}$ for sufficiently large $r$. By the same reasoning, choosing $r$ to be large enough ensures the second inequality. 

    Now, let $r > r_0.$ Then $$\zeta(r) = 1 + 2^{-r} + \sum_{n \geq 3} n^{-r} < 1 + 2 \cdot 2^{-r} < \frac{1+2^{-r}}{1-2^{-r}}$$ and hence $\zeta_{\Q,1}(r) = \zeta(r)(1-2^{-r}) < 1 + 2^{-r}.$ Similarly, we have $$\zeta(r) = 1 + 2^{-r} + 3^{-r} + \sum_{n \geq 4} n^{-r} < 1 + 2^{-r} + 2 \cdot 3^{-r} < \frac{1+3^{-r}}{(1-2^{-r})(1-3^{-r})}$$ and hence $\zeta_{\Q,2}(r) = \zeta(r)(1-2^{-r})(1-3^{-r}) < 1+3^{-r}.$ Consequently, $j_1 \geq 2,$  so $\ell = \max(2,j_1) = j_1.$

    %Note that $j_1 \geq 2$ because $r > r_0.$  The inequality in the definition of $j_1$ fails for $j = 1$ due to the condition on $r$. In addition, the function $a_2(x) := \zeta(x) - \frac{1 + 3^{-x}}{(1 - 2^{-x})(1-3^{-x})}$ also has a unique root greater than 1 (via a similar argument as the one described above); denote it by $r_0^*.$ It is easy to verify that $a_1(x) > a_2(x)$ for $1 < x < 2,$ and estimating $\zeta(x)$ using an integral shows that $r_0 < 2$. Therefore, $r > r_0 > r_0^*,$ so the inequality also fails for $j = 2,$ implying that $j_1 \geq 2,$ so $\ell =j_1.$




    Now, we show that $C_{r,\Q,\chi_{m_{j_1}}} = 2$.  By Lemma \ref{rangeClemma}, we wish to show that $$1 < \log_{p_{j_1}^{-r}} \frac{\zeta_{\Q,j_1}(r)-1}{\zeta_{\Q,j_1}(r) - p_{j_1}^{-r}} < 2.$$ This is equivalent to showing that $\zeta_{\Q,j_1}(r) < 1 + p_{j_1}^{-r}$ and $$\zeta_{\Q,j_1}(r) > \frac{(p_{j_1}^{-r})^2 + p_{j_1}^{-r} + 1}{p_{j_1}^{-r} + 1}.$$ The first inequality follows from the minimality of $j_1$. We proceed to check the second inequality. Let $q = p_{j_1}$ and let $q_* = p_{j_1+1}$. By the definition of $j_1,$ we have $$\prod_{p > q} \frac{1}{1-p^{-r}} > \frac{1 + q_*^{-r}}{1 - q_*^{-r}}.$$ The right hand side of the inequality is a decreasing function in $q_*$, and Bertrand's Postulate tells us that $q_* < 2q,$ so the desired follows from showing $$\frac{1+q^{-r} + (q^{-r})^2}{1 + q^{-r}} < \frac{1 + (2q)^{-r}}{1 - (2q)^{-r}}.$$ Let $x = q^{-r}.$ The inequality above equates to $x^2 + (2-2^r)x + 2 > 0,$ which we need only prove holds for $x \leq 2^{-r}.$ If $1 < r < 2,$ then $2 - 2^r > -2,$ so $x^2 + (2-2^r)x + 2 > x^2 - 2x+ 2 = (x-1)^2 + 1 > 0.$ Now, suppose $r \geq 2.$ Note that $x^2 + (2-2^r)x + 2$ is decreasing on $x \in (0,2^{-r}]$ because its derivative is $2x + 2 - 2^r < 0.$ Therefore, $x^2 + (2-2^r)x + 2 \geq (2^{-r})^2 + (2-2^r)(2^{-r}) + 2 > 0,$ so the inequality holds. 

    Using Lemma \ref{rangeClemma}, we have now shown that for any $r > r_0,$ any integer greater than 1 occurs as $C_{r,\Q,\chi}$ for some character $\chi$. To conclude, we note that there exists a character $\chi$ such that $C_{r,\Q,\chi} = 1.$ Let $$m = \prod_{k \leq j_1} p_k.$$ Because $\chi_m(p_k)=0$ for $k\leq j_1,$ we have $\overline{\sigma_{r,\Q,\chi_m} ( \N)} = \overline{\sigma_{r,\Q,\chi_0}(\mathcal N_{j_1})}.$ As $j_1 = j_{\Q,\chi_0}(r),$ we then apply Lemma \ref{finiteness-first-interval} to conclude $\overline{\sigma_{r,\Q,\chi_m} ( \N)}= [ 1, \zeta_{\Q,j_1}(r)]$ and thus $C_{r,\Q,\chi_m} =1.$
\end{proof}

\begin{comment}
Current thm statement true - real r, using just PRINCIPAL characters
How about real (order-2) characters? 
\end{comment}

\section{A lower bound for $C_{r,\Q,\chi}$} \label{section-lower-1}

In this section, we prove Theorem \ref{lowerbound1}, which gives a lower bound on $C_{r,\Q,\chi}$ for any $r \geq 3$ and any character $\chi$. This result improves upon \cite[Theorem 5.2]{AB19}, which gives a lower bound on $C_{r,\Q}$ for any $r > 1$. Inspired by the proof of \cite[Theorem 5.2]{AB19}, the key mechanism behind the proof of Theorem \ref{lowerbound1} is to partition the natural numbers based on prime factors bounded by functions of $r$ and the modulus of $\chi$, which then capture explicit gaps in the range of the associated divisor function. 
To begin, we prove the following useful lemma. 

\begin{lemma} \label{tail}
    Let $r > 1.$
    \begin{itemize}
        \item[(1)] For any integer $0 < d \leq r-1,$ we have $$d^{-r} > \sum_{n = d+1}^\infty n^{-r}.$$
        \item[(2)] For any odd integer $d \leq 2r-2,$ we have $$d^{-r} > \sum_{n \geq d+2; 2 \nmid n} n^{-r}.$$ 
    \end{itemize}
\end{lemma}

\begin{proof}
    (1) is proved in \cite[Lemma 5.1]{AB19}. (2) follows from noting that $$\sum_{\text{odd } n = d+2}^\infty n^{-r} < \int_{\frac{d-1}{2}}^\infty (2x+1)^{-r} dx = \frac{d^{-r+1}}{2(r-1)},$$ which is less than or equal to $d^{-r}$ when $d \leq 2r-2.$
\end{proof}
We are ready to prove Theorem \ref{lowerbound1}.

\begin{proof}[Proof of Theorem \ref{lowerbound1}]
Let $\theta$ be a nonzero complex number. Then any complex number $x \in \C$ can be written uniquely in the form $a \theta + b \theta i,$ where $a,b \in \R.$ Define $\Re_\theta$ by $\Re_\theta(a\theta + b \theta i) = a.$


We begin by partitioning $\N$. To do so, we first define the map $B_1: \N \rightarrow \mathcal{P}(\N)$ by $B_1(x) = \{d \in \N : d \mid x, d \leq r-1, \gcd(d,m)=1\}.$ Next, we define the map $B_2 : \N \rightarrow \mathcal {P} (\N)$ by $$B_2(n) = \begin{cases}
    \{d \in \N : d \mid n, r-1 \leq d \leq 2r-2, \gcd(d, 2m) = 1\} & \text{ if } 2 \not\in B_1(n) \\
    \emptyset &\text{ if }2 \in B_1(n).
\end{cases}$$
Finally, define $A(n) = (B_1(n), B_2(n)).$
Note that this determines a partition of $\N$ whose parts are $A^{-1}(S)$ for $S \in A(\N).$ We wish to show that for any distinct $S_1, S_2 \in A(\N),$ the sets $\sigma_{r,\Q,\chi}(A^{-1}(S_1))$ and $\sigma_{r,\Q,\chi}(A^{-1}(S_2))$ have disjoint closures. Write $S_1 = A(x)$ and $S_2 = A(y)$ for distinct $x,y \in \N.$ 

    We consider separately the cases $B_1(x) \neq B_1(y)$ and $B_1(x) = B_1(y)$. Assume first that $B_1(x) \neq B_1(y)$.  Suppose without loss of generality that the smallest integer $d_0 \leq r-1$ coprime to $m$ that divides precisely one of $x$ and $y$ lies in $B_1(x)$. Let $\theta = \chi(d_0).$ Take $x'$ such that $B_1(x') = B_1(x)$ and $y'$ such that $B_1(y') = B_1(y).$ For each $d < d_0,$ the difference of the terms for $d$ in $\sigma_{r,\Q,\chi} (x') - \sigma_{r, \Q,\chi}(y')$ is 0: if $\gcd(d,m) > 1,$ then $\chi(d) = 0,$ and otherwise, $d \in B_1(x')$ if and only if $d \in B_1(y')$ by the minimality of $d_0.$ Then we have $$\sigma_{r,\Q,\chi} (x') - \sigma_{r, \Q,\chi}(y') = \theta d_0^{-r} + \sum_{d \mid x'; d > d_0} \chi(d)  d^{-r} - \sum_{d \mid y'; d > d_0} \chi(d) d^{-r}.$$ Applying $\Re_\theta$ to both sides, we obtain $$\Re_\theta(\sigma_{r,\Q,\chi}(x')) - \Re_\theta(\sigma_{r,\Q,\chi}(y')) \geq d_0^{-r} - \sum_{d > d_0} d^{-r} > 0$$ by Lemma \ref{tail}. 
Thus, the closures of $\sigma_{r,\Q,\chi}(A^{-1}(A(x))$ and $\sigma_{r,\Q,\chi}(A^{-1}(A(y))$ are disjoint. 

    It remains to consider the case $B_1(x) = B_1(y).$ Then $B_2(x) \neq B_2(y),$ so one of these sets is nonempty, meaning $2 \not\in B_1(x) = B_1(y).$ If $m$ is even, then $\chi(d) = 0$ for even $d$. If $m$ is odd, because  $r-1 \geq 2$, the fact that $2 \not\in B_1(x) = B_1(y)$ implies that $x$ and $y$ have no even divisors.  The argument then proceeds in a similar fashion to the case $B_1(x) \neq B_1(y).$ Let $d_0$ be the smallest integer that is in precisely one of the two sets $B_2(x)$ and $B_2(y).$ Without loss of generality, we assume $d_0 \in B_2(x) \setminus B_2(y).$ Let $\theta = \chi(d_0),$ and take $x'$ such that $A(x') = A(x)$ and $y'$ such that $A(y') = A(y).$ When comparing $\Re_\theta(\sigma_{r,\Q,\chi}(x'))$ and $\Re_\theta(\sigma_{r,\Q,\chi}(y'))$, we use the fact that the contributions from divisors $d < d_0$ are the same. In addition, there are no contributions from even divisors. Indeed, if $m$ is even, then $\chi(d) = 0$ for even $d$; if $m$ is odd, because  $r-1 \geq 2$, the fact that $2 \not\in B_1(x) = B_1(y)$ implies that $x$ and $y$ have no even divisors.  More explicitly, we compute 
    \begin{align*}
        \Re_\theta(\sigma_{r,\Q,\chi}(x'))-\Re_\theta(\sigma_{r,\Q,\chi}(y')) &= \Re_\theta \left( \sum_{d \mid x'; d < d_0} d^{-r} \chi(d) + d_0^{-r} \theta + \sum_{d \mid x'; d \geq d_0+2; d \text{ odd}} d^{-r}\chi(d)\right) \\
        &\text{ \; } - \Re_\theta \left( \sum_{d \mid y'; d < d_0} d^{-r} \chi(d) +  \sum_{d \mid y'; d \geq d_0+2; d \text{ odd}} d^{-r}\chi(d)\right) \\
        &= d_0^{-r} + \Re_\theta\left( \sum_{d \mid x'; d \geq d_0+2; d \text{ odd}} d^{-r}\chi(d) - \sum_{d \mid y'; d \geq d_0+2; d \text{ odd}} d^{-r}\chi(d)\right) \\
        &\geq d_0^{-r} - \sum_{d \geq d_0 + 2 ; d \text{ odd}} d^{-r},
    \end{align*}
    which is positive by Lemma \ref{tail}. Thus, we have $C_{r, \Q, \chi} \geq |A(\mathbb N)|.$ 

    It remains to give a lower bound for $|A(\N)|.$ Let $P$ be the set of odd primes $p \leq 2r-2$ that are coprime to $m$. For each subset $S \subseteq P,$ define $$n_S = \prod_{p \in S} p .$$ Because $n_S$ is odd, $2 \not\in B_1(n_S),$ so $B_2(n_S) = \{d \in \N : d \mid n_S, r-1 \leq d \leq 2r-2, \gcd(d, 2m) = 1\}.$ Each prime $p \in S$ is recorded either in $B_1(n_S)$ (if $p \leq r-1$) or in $B_2(n_S)$ (if $r-1\leq p \leq 2r-2$). Thus, distinct choices of subsets $S$ of $P$ give rise to distinct values of $A(n_S)$. Consequently, we have $$C_{r, \Q, \chi} \geq |A(\mathbb N)| \geq 2^{|P|} = 2^{\pi_m(2r-2) - \delta_m},$$ where $\delta_m$ equals 0 if $m$ is even and equals 1 if $m$ is odd.
\end{proof}

\section{A lower bound for $C_{r,K}$} \label{section-lower-2}



In this section, we prove Theorem \ref{lowerbound2}, which gives a lower bound for $C_{r,K}$ for finite Galois extensions $K$ of $\Q$ and generalizes the lower bound that \cite[Theorem 4.1]{AB19} gives when $K = \Q$. The proof's key mechanism is inspired by that of the previous section and of \cite[Theorem 4.1]{AB19}. We partition the set of integral ideals in \(\mathcal O_K\)
based on the multiplicities of their divisor ideals whose norms are bounded by a
function \(h_K(r)\) depending on \(K\) and \(r\). This partition captures explicit
gaps in the range of the associated divisor function; thus, the problem reduces to
lower-bounding \(h_K(r)\) and counting how many distinct partition classes are
created.
The new feature when handling a number field \(K\), rather than just \(\mathbb Q\),
is that distinct ideals can have the same norm. Then in order to bound
\(h_K(r)\), we control the number of ideals of each fixed norm using the
coefficients of the Dedekind zeta function. Finally, we use Chebotarev's density
theorem to estimate the number of rational primes that split completely in \(K\), giving a 
lower bound for the number of gaps created.

 We begin by defining some notation that we use throughout this section. We fix a degree-$s$ Galois extension $K/ \Q. $ 
For $n \geq 1,$ we define $a_K(n)$ to be the number of integral ideals $I$ in $K$ such that $N(I) = n.$ Then the Dedekind zeta function can be written as $$\zeta_K(r) = \sum_{n=1}^\infty \frac{a_K(n)}{n^r}.$$ In addition, we define $b_K(n)$ be the number of prime ideals $\mathfrak p$ in $K$ such that $N(\mathfrak p) = n.$ Lastly, for any ideal $I \subseteq \mathcal O_K,$ we define $m_I(n)$ to be the number of integral ideals $J$ in $K$ for which $J \mid I$ and $N(J) = n$. Then the divisor function can be written as $$\sigma_{r,K}(I) = \sum_{n=1}^\infty \frac{m_I(n)}{n^r}.$$

We first prove a general ``theoretical" lower bound on $C_{r,K}$ that generalizes the partition method used in the previous section. 

\begin{definition}
    Fix $r > 1.$ We define $h_K(r)$ to be the largest integer $h \geq 1$ such that for each integer $2 \leq d \leq h,$ the inequality $$d^{-r} > \sum_{n > d} \frac{a_K(n)}{n^r}$$ holds. If no such $h \geq 2$ exists, then define $h_K(r) = 1.$
\end{definition}

\begin{theorem} \label{partition}
    For any $r > 1,$ we have $$C_{r,K} \geq \prod_{n=2}^{h_K(r)} (b_K(n) + 1).$$
\end{theorem}

\begin{proof}
    Let $h = h_K(r).$ Define the map $A : I_K \ra \Z_{\geq 0}^{h-1}$ by $$A(I) = (m_I(2), m_I(3), \ldots, m_I(h)).$$ Note that this determines a partition of $I_K$ whose parts are $A^{-1}(S)$ for $S \in A(I_K).$ We wish to show that for any distinct $S_1, S_2 \in A(I_K),$ the sets $\sigma_{r,K}(A^{-1}(S_1))$ and $\sigma_{r,K}(A^{-1}(S_2))$ have disjoint closures. Write $S_1 = A(\mathfrak a)$ and $S_2 = A(\mathfrak b)$ for distinct $\mathfrak a, \mathfrak b \in I_K.$ Let $d_0$ be the smallest integer $2 \leq d_0 \leq h$ such that $m_\fa(d_0) \neq m_\bb(d_0),$ and we assume without loss of generality that $m_\fa(d_0) > m_\bb(d_0).$ 

    Now, let $\fa'$ and $\bb'$ be ideals in $K$ such that $A(\fa') = A(\fa)$ and $A(\bb') = A(\bb).$ From the definition of $d_0,$ we have $m_{\fa'}(d_0) - m_{\bb'}(d_0) \geq 1.$ Thus, we observe 
    \begin{align*}
        \sigma_{r,K}(\fa') - \sigma_{r,K}(\bb') &= \sum_{n=1}^\infty \frac{m_{\fa'}(n) - m_{\bb'}(n)}{n^r} \\
        &\geq d_0^{-r} - \sum_{n = d_0+1}^\infty \frac{m_{\bb'}(n)}{n^r} \\
        &\geq d_0^{-r} - \sum_{n = d_0+1}^\infty \frac{a_K(n)}{n^r}.
    \end{align*} Because $d_0 \leq h_K(r),$ the definition of $h_K(r)$ implies that the last expression above is positive and thus that $\sigma_{r,K}(\fa') > \sigma_{r,K}(\bb').$

    Thus, we have shown that the closures of $\sigma_{r,K}(A^{-1}(S_1))$ and $\sigma_{r,K}(A^{-1}(S_2))$ are disjoint. It follows that $C_{r,K} \geq |A(I_K)|,$ so it remains to determine a lower bound for $|A(I_K)|.$ We do so in a fashion analogous to that of the previous section. For each $2 \leq n \leq h,$ denote the prime ideals of norm $n$ by $$\mathfrak{p}_{n, 1}, \ldots, \mathfrak{p}_{n, b_K(n)}.$$ For each tuple $(i_n)_{2 \leq n \leq h},$ where $0 \leq i_n \leq b_K(n),$ define the ideal $$I_{(i_n)} = \prod_{n=2}^h \prod_{j=1}^{i_n} \mathfrak{p}_{n, j}.$$ We check that distinct tuples $(i_n)$ give distinct values of $A(I_{(i_n)}).$ Let $\textbf{i} =(i_n)$ and $\textbf{j} = (j_n)$ be distinct tuples, and let $n_0$ be the smallest integer for which $i_{n_0} \neq j_{n_0}.$ We show $m_{I_\textbf{i}}(n_0) \neq m_{I_\textbf{j}}(n_0).$ The prime ideals dividing $I_\textbf{i}$ and $I_\textbf{j}$ that are of norm less than $n_0$ are identical. Thus, the ideals dividing $I_\textbf{i}$ and $I_\textbf{j}$ that are not prime and are of norm equal to $n_0$ are also identical. Hence, $m_{I_\textbf{i}}(n_0) - m_{I_\textbf{j}}(n_0)$ equals the difference in the number of prime ideals dividing $I_\textbf{i}$ and $I_\textbf{j}$ that have norm $n_0,$ which is $i_{n_0} - j_{n_0} \neq 0.$

    
    Therefore, $$C_{r,K} \geq \#\{(i_n)_{2 \leq n \leq h} : 0 \leq i_n \leq b_K(n)\} = \prod_{n=2}^h (b_K(n) + 1).$$\end{proof}

Now, we aim to bound $a_K(n)$ in terms of $n$.
To do this, we first define the function $$d_s(n) = \#\{(n_1, \ldots, n_s) \in \mathbb N^s : n_1 n_2 \cdots n_s = n\}$$ for any positive integer $s. $ The function $d_s(n)$ is clearly multiplicative. For any prime $p$ and nonnegative integer $e$, we see that $$d_s(p^e) = \# \{(e_1, \ldots, e_s) \in \Z_{\geq 0} : e_1 + \dots + e_s = e\} = \binom{e + s-1}{s-1}.$$ Therefore, we have the general formula $$d_s(n) = \prod_{p^{e_p} \|n} \binom{e_p+s-1}{s-1}.$$ 

\begin{lemma}
    Let $K/\Q$ be a degree-$s$ Galois field extension. For any $n \geq 1,$ we have $a_K(n) \leq d_s(n).$
\end{lemma}

\begin{proof}
    The Dedekind zeta function can be expressed as $$\zeta_K(x) = \prod_{\mathfrak p \subseteq \mathcal O_K} (1 - N(\mathfrak p)^{-x})^{-1}.$$ Because $a_K$ and $d_s$ are multiplicative, we fix a rational prime $p$ and prove $a_K(p^a) \leq d_s(p^a)$ for any positive integer $a$. Let $f_p$ be the inertia degree of a prime ideal in $K$ lying above $p$, and let $r_p$ be the number of distinct prime ideals in $K$ lying above $p$. Then $$\zeta_K(x) = \prod_p  (1 - p^{-f_p x})^{-r_p}.$$ 
    Note that $a_K(p^a)$ is precisely the coefficient of $p^{-ax}$ in $\zeta_K(x)$ and is thus the coefficient of $p^{-ax}$ in the $p$-Euler factor $(1 - p^{-f_p x})^{-r_p}.$ Because $f_p \geq 1$ and $r_p \leq s,$ the coefficient of $p^{-ax}$ in the expansion of $(1 - p^{-f_p x})^{-r_p}$ is less than or equal the coefficient of $p^{-ax}$ in the expansion of $(1 - p^{-x})^{-s}.$ The latter coefficient is precisely $d_s(p^a),$ as desired.
\end{proof}

Next, we bound $d_s(n)$ in terms of $n$. 

\begin{lemma}
    Fix a positive integer $s$. For any $\varepsilon > 0,$ there exists a constant $\eta_{s, \varepsilon} > 0$ such that $d_s(n) \leq \eta_{s, \varepsilon} n^\varepsilon$ for any $n \geq 1.$
\end{lemma}

\begin{proof}
    Fix $\varepsilon > 0.$ Recall the closed formula $$d_s(n) = \prod_{p^{e_p} \|n} \binom{e_p+s-1}{s-1}.$$ Each binomial coefficient $\binom{e+s-1}{s-1}$ is bounded above by $e^{s-1}$ for large enough $e$. This polynomial growth is dominated eventually by exponential growth, so there exists $E$ such that if $e > E,$ then $$\binom{e+s-1}{s-1} \leq 2^{\frac{e \varepsilon}{2}}.$$ We then split the product in $d_s(n)$ using $E$; that is, we write $d_s(n) = d_{s,1}(n) d_{s,2}(n),$ where we define $$d_{s,1}(n) = \prod_{p^{e_p} \|n, 1 \leq e_p < E} \binom{e_p+s-1}{s-1} \text{ and }d_{s,2}(n) = \prod_{p^{e_p} \|n, e_p \geq E} \binom{e_p+s-1}{s-1}.$$ We can bound the second product as follows. $$d_{s,2}(n) \leq \prod_{p^{e_p} \| n} 2^\frac{e_p \varepsilon}{2} \leq \prod_{p^{e_p} \| n} p^\frac{e_p \varepsilon}{2} = n^\frac{\varepsilon}{2}$$ We proceed to bound the first product. Let $$M = \max_{1 \leq e < E} \binom{e+s-1}{s-1}.$$ Then $d_{s,1}(n) \leq M^{\omega(n)},$ where $\omega(n)$ is the number of distinct prime factors of $n$. Using the standard upper bound $\omega(n) = O(\frac{\log n}{\log \log n}),$ we find that 
    \begin{align*}
        d_{s,1}(n) & \leq \exp(\log M \cdot \omega(n)) \\
        &= \exp\left( O_{s, \varepsilon} \left( \frac{\log n}{\log \log n}\right)\right) \\
        &\leq \exp\left( \delta_{s, \varepsilon} \frac{\log n}{\log \log n}\right) \text{ for some constant $\delta_{s,\varepsilon}$ and large enough }n \\
        &\leq \exp\left( \frac{\varepsilon}{2} \log n \right)\text{ for large enough }n.
    \end{align*} Therefore, we have $d_s(n)  = d_{s,1}(n) d_{s,2}(n) \leq n^\varepsilon$ for large enough $n$. Consequently, accounting for the finitely many smaller values of $n$, there exists a constant $\eta_{s, \varepsilon} > 0$ such that $d_s(n) \leq \eta_{s, \varepsilon} n^\varepsilon$ for any $n \geq 1.$
\end{proof}

Thus, we have shown that for any finite Galois extension $K/\Q$ and $\varepsilon > 0,$ there exists a constant $\eta_{K, \varepsilon} := \eta_{s,\varepsilon}$ such that $a_K(n) \leq \eta_{K, \varepsilon} n^\varepsilon$ for any positive integer $n$. We use this bound to show that the function $h_{K, \varepsilon}(r)$ we define next is a lower bound for $h_K(r)$. Then we apply Theorem \ref{partition} and Chebotarev's density theorem  to prove Theorem \ref{lowerbound2}. 

\begin{definition}
    Let $K/\Q$ be a finite Galois extension, and let $\varepsilon > 0.$ For $r > \varepsilon + 1,$ define $$h_{K, \varepsilon}(r) = \left\lfloor \left( \frac{r-\varepsilon-1}{\eta_{K,\varepsilon}} \right)^{\frac{1}{1+\varepsilon}} \right\rfloor.$$
\end{definition}

\begin{proof}[Proof of Theorem \ref{lowerbound2}]
    We show $h_{K, \varepsilon}(r) \leq h_K(r).$ Let $2 \leq d \leq h_{K, \varepsilon}(r).$ Because $a_K(n) \leq \eta_{K, \varepsilon}n^\varepsilon$ for any positive integer $n$, we have \begin{align*}
        \sum_{n > d} \frac{a_K(n)}{n^r} &\leq \eta_{K, \varepsilon} \sum_{n > d} n^{-r + \varepsilon} \\
        &< \eta_{K, \varepsilon}\int_d^\infty x^{-r+\varepsilon} \; dx \\
        &= \frac{\eta_{K, \varepsilon}}{r-\varepsilon-1} d^{-r+\varepsilon+1},
    \end{align*} where the last equality holds because $r > \varepsilon + 1.$ The fact that $d \leq h_{K, \varepsilon}(r)$ implies $$\frac{\eta_{K, \varepsilon}}{r-\varepsilon-1} d^{-r + \varepsilon + 1} \leq d^{-r}.$$ The computation above thus gives $$\sum_{n > d} \frac{a_K(n)}{n^r} < d^{-r},$$ verifying $h_{K, \varepsilon}(r) \leq h_K(r).$

    Recall that by Theorem \ref{partition}, we have $$C_{r,K} \geq \prod_{n=2}^{h_{K,\varepsilon}(r)} (b_K(n) + 1).$$ We then restrict this product to the rational primes that split completely in $K$ in order to apply Chebotarev's density theorem. In particular, we have \begin{align*}
        C_{r,K} &\geq \prod_{p \leq h_{K, \varepsilon}(r) \text{ splits completely in }K} (b_K(p) + 1) \\
        &= (s+1)^{\#\{ p \leq h_{K, \varepsilon}(r) : p \text{ splits completely in }K\}}.
    \end{align*} Notice  that $h_{K, \varepsilon}(r) \asymp_{K,\varepsilon} r^{1/ (1 + \varepsilon)}$; that is, $h_{K,\varepsilon}(r)$ is bounded above and below by $(K,\varepsilon)$-dependent constant multiples of $r^{1/(1+\varepsilon)}$ for $r$ sufficiently large.
    Chebotarev's density theorem then gives us 
    \begin{align*}
        \#\{ p \leq h_{K, \varepsilon}(r) : p \text{ splits completely in }K\} &\sim \frac{1}{s} \frac{h_{K,\varepsilon}(r)}{\log h_{K, \varepsilon}(r)} \\
        &\asymp_{K,\varepsilon} \frac{r^{1/(1+ \varepsilon)}}{\log r}.
    \end{align*} Because $s$ is dependent on $K$, we conclude that there exist constants $\mu_{K,\varepsilon}$ and $r_{K,\varepsilon}$ for which $$\log C_{r,K} > \mu_{K,\varepsilon} \left( \frac{r^{1/(1+ \varepsilon)}}{\log r} \right)$$ holds for any $r > r_{K,\varepsilon}.$
\end{proof}


\section{Unboundedness of $C_{r,K}$ when $r$ and $[K:\Q]$ are fixed}\label{section-last}




The previous section showed that for a fixed number field $K$, the number of connected components $C_{r,K}$ tends to $\infty$ as $r$ tends to $\infty.$ A natural next question is whether $C_{r,K}$ exhibits the same behavior when $r$ is fixed and $K$ is varied. In this section, we prove Theorem \ref{fix-s}, which answers this question in the positive. Indeed, we show the stronger result that $C_{r,K}$ can be arbitrarily large when $r$ and $[K:\Q]$ are fixed and $K$ is varied.
To prove Theorem \ref{fix-s}, we define a notion of \textit{mighty norms} that generalizes terminology defined in \cite{Z18}, after which we outline the main ideas of the proof.  

\begin{definition}
    Let $d > 1$ be the norm of some prime ideal in $K$. We say $d$ is an \textit{$r$-mighty norm} of $K$ if \begin{equation} \label{mighty}
        1 + d^{-r} > \prod_{N(\mathfrak p) > d} (1 - N(\mathfrak p)^{-r})^{-1}.
    \end{equation}
\end{definition}

Like Zubrilina's definition, the existence of $r$-mighty norms correspond to the existence of gaps in $\sigma_{r,K}(I_K),$ which we make precise below. 

\begin{lemma}\label{lemma-gaps}
    Let $d_1 < d_2 < \dots < d_M$ be $r$-mighty norms of $K$. Then $C_{r, K} \geq M+1.$
\end{lemma}

\begin{proof}
    Fix $d := d_i$ for some $1 \leq i \leq M.$ We first show the existence of a gap in $\overline{\sigma_{r,K}(I_K)}$ associated to $d$. Let $S$ be the set of ideals $I$ in $K$ such that some prime ideal factor $\mathfrak q$ of $I$ satisfies $N(\mathfrak q) \leq d.$ For any $I \in S,$ we have $\sigma_{r,K}(I) \geq 1+d^{-r}.$ Meanwhile, for any $I \in I_K \setminus S,$ we have $\sigma_{r,K}(I) \leq \prod_{N(\mathfrak p) > d} (1 - N(\mathfrak p)^{-r})^{-1}.$ Therefore, $\overline{\sigma_{r,K}(S)} \cap \overline{\sigma_{r,K}(I_K \setminus S)} = \emptyset.$ 

    Now, we verify that for any $1 \leq i < j \leq M,$ the gap intervals $$\left( \prod_{N(\mathfrak p) > d_j} (1 - N(\mathfrak p)^{-r})^{-1}, 1 + d_j^{-r} \right) \text{\; and \;} \left( \prod_{N(\mathfrak p) > d_i} (1 - N(\mathfrak p)^{-r})^{-1}, 1 + d_i^{-r} \right)$$ are disjoint. We see easily that $$\prod_{N(\mathfrak p) > d_i} (1 - N(\mathfrak p)^{-r})^{-1} \geq (1 - d_j^{-r})^{-1} > 1 + d_j^{-r},$$ so the intervals are disjoint. Thus, there are at least $M$ many gap intervals in $\overline{\sigma_{r,K}(I_K)},$ so $C_{r,K} \geq M+1.$
\end{proof}


For an arbitrary positive integer $M$, we aim to construct a sequence of rational primes $p_1 < p_2 < \cdots < p_M$ and a degree-$s$ number field $K$ in which $p_1,\ldots,p_M$ are $r$-mighty norms. By Lemma \ref{lemma-gaps}, this would immediately give Theorem \ref{fix-s}.
The idea is to exploit the effect of a prime ideal's splitting behavior on its norm. Notice that rational primes $p$ that split completely in a number field $K/\Q$ have ``smaller" norms $p$, whereas rational primes that are inert in $K$ have ``larger" norms $p^s$. Hence, if the primes $p_1 < p_2 < \dots < p_M$ are spaced widely enough apart and if they split completely while nearby primes remain inert in $K$, then $p_1,p_2,\ldots,p_M$ are $r$-mighty norms of $K$. The first lemma we introduce achieves the first goal of constructing such a sequence of rational primes, and the second lemma we introduce achieves the second goal of showing the existence of a number field $K$ with the desired splitting behaviors. 

\begin{lemma} \label{technical}
    Fix a real number $r > 1,$ an integer $s \geq 2,$ and an integer $M \geq 1.$ There exists a sequence of rational primes $p_1 < p_2 < \dots < p_M$ satisfying the following conditions. 
    \begin{enumerate}
        \item \begin{enumerate}
            \item $p_1 > \max(2,s)^s.$
            \item $\log(1-x)^{-1} \leq 2x$ for $x = p_1^{-r}.$
            \item $\log(1+x) \geq \frac{x}{2}$ for $x = p_1^{-r}.$
        \end{enumerate} \noindent (Note that (b) and (c) thus also hold for any $0 \leq x \leq p_1^{-r}.$)
        \item For each $2 \leq i \leq M,$ we have $p_i^{1/s} > p_{i-1}.$ In particular, defining for each $1 \leq i \leq M$ the set $S_i = \{q \text{ prime} : p_i^{1/s} < q \leq p_i\},$ the sets $S_1,\ldots,S_M$ are pairwise disjoint.
        \item For any $1 \leq j < i \leq M,$ we have $$s \cdot \sum_{q \in S_i} q^{-r} < \frac{1}{12M} p_j^{-r}.$$
        \item For any $1 \leq i \leq M,$ we have $$\sum_{\text{prime }q > p_i} q^{-sr} < \frac{1}{12}p_i^{-r}.$$
    \end{enumerate}
\end{lemma}

Before we give a technical proof of this lemma, we briefly explain the later necessity of each condition. Recall that our goal is to show $p_1, p_2, \ldots, p_M$ are $r$-mighty norms for some degree-$s$ number field $K$.
Condition (1) allows us to compare the right-hand side of Equation (\ref{mighty}) for $d = p_i$ with simpler expressions in $p_i,$ such as $\sum_{N(\mathfrak p) > p_i} N(\mathfrak p)^{-r}.$ Then conditions (2), (3), and (4) allow us to give an upper bound on $\sum_{N(\mathfrak p) > p_i} N(\mathfrak p)^{-r}$ that we then compare to the left-hand side of Equation (\ref{mighty}) for $d = p_i.$ 

\begin{proof}[Proof of Lemma \ref{technical}]
    We construct the desired sequence of primes recursively. We first construct the desired $p_1.$ Condition (1) can be imposed by taking $p_1$ sufficiently large, because $$\lim_{x \ra 0} \frac{\log(1-x)^{-1}}{x} = \lim_{x \ra 0} \frac{\log(1+x)}{x} = 1.$$
    
    Now, suppose we have chosen $p_1, \ldots, p_{i-1}$ to satisfy the above conditions; we aim to choose $p_i$ such that $p_1,\ldots,p_{i}$ satisfy the four conditions. Condition (2) clearly just requires $p_i$ to be sufficiently large. Next, we consider condition (3). Because $r > 1,$ we may write $$\sum_{q \in S_i} q^{-r} \leq \sum_{q  > p_i^{1/s}} q^{-r} < \int_{p_i^{1/s}}^\infty x^{-r} \;dx = \frac{(p_i^{1/s})^{-r+1}}{r-1} \ll_r p_i^{\frac{1-r}{s}}.$$ Therefore, we have $$s \cdot \sum_{q \in S_i} q^{-r}  \ll_{r, s} p_i^{\frac{1-r}{s}}.$$ The right-hand side tends to 0 as $p_i$ tends to $\infty,$ so we can choose $p_i$ sufficiently large to bound it from above by $\frac{1}{12M} \min_{1 \leq j < i} p_j^{-r}.$ Lastly, we consider condition (4). Similarly, because $s \geq 2$ and $r > 1$ and thus $sr-1 > 1,$ we see that $$\sum_{\text{prime }q > p_i} q^{-sr} \ll_r p_i^{-sr+1}.$$ Because $sr-1 > r,$ we can choose $p_i$ sufficiently large to bound the left-hand side from above by $\frac{1}{12}p_i^{-r}.$
\end{proof}

\begin{lemma} \label{ntlemma}
    Let $s \geq 2$ be an integer, and let $S$ and $T$ be disjoint finite sets of rational primes. Furthermore, assume each prime in $S$ is greater than $s$. Then there exists a degree-$s$ extension $K/\Q$ such that each prime in $S$ splits completely in $K$ and each prime in $T$ remains inert in $K$. 
\end{lemma}
\begin{proof}
    We aim to construct a monic irreducible polynomial  $f(x) \in \Z[x]$ such that, modulo each prime in $S \cup T,$ its factorization type corresponds to the desired splitting behavior. We then can apply the Dedekind--Kummer theorem to show that $\Q(\alpha)$, where $\alpha$ is a root of $f(x),$ is a desired choice of $K$. 


    For each $p \in S,$ we choose a monic degree-$s$ polynomial $f_p(x) = x^s + c_{p,s-1}x^{s-1} + \cdots + c_{p,0} \in \mathbb F_p[x]$ that splits into $s$ distinct linear factors over $\mathbb F_p.$ This is feasible because $s < p.$ Furthermore, for each $p \in T,$ we choose a monic irreducible polynomial $f_p(x)  = x^s + c_{p,s-1} x^{s-1} + \cdots + c_{p,0} \in \mathbb F_p[x]$ of degree $s$. Finally, fix some prime $q \not\in S \cup T,$ and let $f_q(x) = x^s + q.$ We show there exists a monic degree-$s$ polynomial $f(x) \in \Z[x]$ such that 
    $$f(x) \equiv \begin{cases}
        f_p(x) \pmod{p}& \text{ for } p \in S \cup T \\
        f_q(x) \pmod{q^2}
    \end{cases}$$
    For each $0 \leq i \leq s-1,$ we apply the Chinese remainder theorem to obtain $n_i \in \Z$ such that $n_i \equiv c_{p,i} \pmod{p}$ for each $p \in S \cup T$ and $n_i \equiv q \pmod{q^2}$ if $i = 0$ and $n_i \equiv 0 \pmod{q^2}$ else. Then $f(x) = x^s + n_{s-1}x^{s-1} + \cdots + n_0 \in \Z[x]$ satisfies the desired congruences.  
    
    Note that $f$ is an Eisenstein polynomial because $f(x) \equiv x^s + q \pmod{q^2},$ so it is irreducible. Let $\alpha$ be a root of $f,$ and let $K  = \Q(\alpha).$ To apply the Dedekind--Kummer theorem to $K$ for each $p \in S \cup T$, we must check that $p \nmid [\mathcal O_K : \Z[\alpha]].$ Fix $p \in S \cup T$. We have the identity $$\text{disc}(f) = [\mathcal O_K : \Z[\alpha]]^2 \cdot \text{disc}(K).$$ Now, notice that $f(x) \pmod{p}$ is squarefree; when $p \in S,$ this is by construction, and when $p \in T,$ we note that any irreducible polynomial over a finite field is separable and thus squarefree. Therefore, $p \nmid \text{disc}(f)$. Using the fact that $\text{disc}(K) \in \Z$ because $\alpha$ is an algebraic integer, we obtain $p \nmid [\mathcal O_K : \Z[\alpha]].$ Thus, the Dedekind--Kummer theorem applies at each $p \in S \cup T.$ If $p \in S,$ then $f(x) \equiv f_p(x) \pmod{p}$ splits into $s$ distinct linear factors, so $p$ splits completely in $K$; if $p \in T,$ then $f(x) \equiv f_p(x) \pmod{p}$ is irreducible, so $p$ remains inert in $K$. 
\end{proof}

We are ready to prove Theorem \ref{fix-s}.

\begin{proof}[Proof of Theorem \ref{fix-s}]
    Let $M$ be an arbitrary positive integer. By Lemma \ref{technical}, we have a sequence of rational primes $p_1 < p_2 < \dots < p_M$ satisfying the four conditions specified in the lemma's statement. Now, note that we can choose a real number $X > p_M$ such that $$s \cdot \sum_{q > X} q^{-r} < \frac{1}{12} p_M^{-r}.$$ The left-hand side is bounded above by some $r$-dependent constant multiple of $X^{-r+1}$ and thus tends to 0 as $X$ tends to $\infty,$ so making $X$ large enough ensures the validity of the above inequality. 
    Now, let $$S = \bigcup_{i=1}^M S_i,$$ where $S_i$ are defined in the statement of Lemma \ref{technical}, and let $T$ be the set of rational primes $q \not\in S$ such that $q \leq X.$ Applying Lemma \ref{ntlemma} to these choices of $S$ and $T$, we have a degree-$s$ extension $K/\Q$ such that $p_1, \ldots, p_M$ split completely in $K$ and any rational prime $q \not\in S$ for which $q \leq X$ is inert.  

    Fix $1 \leq i \leq M.$ Because $p_i$ splits completely in $K$, it is a norm of a prime ideal in $K$. We prove that $p_i$ is an $r$-mighty norm of $K$; this would imply $C_{r,K} \geq M+1,$ proving the theorem. Recall that this requires showing $$\prod_{N(\mathfrak p) > p_i} (1 - N(\mathfrak p)^{-r})^{-1} < 1 + p_i^{-r}.$$ By condition (1) in Lemma \ref{technical}, we have $\log(1-x)^{-1} \leq 2x$ for any $x \leq p_1^{-r}.$ Therefore, it suffices to show \begin{equation} \label{eq-suffices}
        2 \sum_{N(\mathfrak p) > p_i} N(\mathfrak p)^{-r} < \log(1 + p_i^{-r}).
    \end{equation} We proceed to bound the sum on the left-hand side by considering the three following types of rational primes $q$ that $\mathfrak p$ can lie over. 
    \begin{itemize}
        \item $q \leq X$ is a rational prime that splits completely in $K$. The norm of any prime ideal lying above $q$ is $q$, so by condition (2) of Lemma \ref{technical}, if $q > p_i,$ then $q \in S_j$ for some $j > i$. Therefore, the total contribution of prime ideals $\mathfrak p$ lying above such $q$ satisfying $N(\mathfrak p) > p_i$ to the summation is bounded from above by $$s \cdot \sum_{j > i} \sum_{q \in S_j} q^{-r} < \sum_{j > i} \frac{1}{12M} p_i^{-r} < \frac{1}{12}p_i^{-r},$$ where the first inequality is condition (3) of Lemma \ref{technical}.
        \item $q \leq X$ is a rational prime that remains inert in $K$. If $q \leq p_i^{1/s},$ then $N(q\mathcal O_K) = q^s \leq p_i,$ so $q \mathcal O_K$ does not contribute to the summation. Because $q \leq X$ is inert in $K$, it does not lie in $S_i,$ so $q \not\in (p_i^{1/s}, p_i].$ Lastly, if $q > p_i$, then the sole prime ideal $q \mathcal O_K$ lying above $q$ contributes $q^{-sr}$ to the summation. The total contribution of these inert $q \leq X$ is $$\sum_{\text{prime }q > p_i} q^{-sr} < \frac{1}{12} p_i^{-r},$$ where the inequality is condition (4) of Lemma \ref{technical}.
        \item $q > X$ is a rational prime. Then $q > p_i,$ so any prime ideal $\mathfrak p$ lying above $q$ satisfies $N(\mathfrak p) > p_i.$ There are at most $s$ prime ideals in $K$ lying above $q$, and each one has norm at least $q$. Thus, the contribution of such ideals to the summation is bounded above by $$s \cdot \sum_{q > X} q^{-r} < \frac{1}{12} p_M^{-r} \leq \frac{1}{12} p_i^{-r},$$ where the middle inequality follows from the definition of $X$.  
    \end{itemize}
    Summing the contributions from these three cases, we conclude that $$\sum_{N(\mathfrak p) > p_i} N(\mathfrak p)^{-r} < \frac{1}{4} p_i^{-r}.$$ Then Equation (\ref{eq-suffices}) holds if $$2 \left(\frac{1}{4} p_i^{-r}\right) = \frac{1}{2} p_i^{-r} < \log(1 + p_i^{-r}).$$ This holds true by condition (1) in Lemma \ref{technical}, so we have shown $p_i$ is an $r$-mighty norm of $K$. 
\end{proof}


\begin{comment}
\begin{lemma} \label{ideal1.75}
    Let $r \geq 2.$ For any $d \leq \sqrt{r-2},$ we have $$d^{-r} > \sum_{n=d+1}^\infty n^{-r+1}.$$
\end{lemma}

\begin{proof}
    We may bound the right-hand side by  $$\sum_{n=d+1}^\infty n^{-r+1} < \int_{d}^\infty x^{-r+1} dx = \frac{d^{-r+2}}{r-2},$$ which is less than or equal to $d^{-r}$ precisely when  $d \leq \sqrt{r-2}.$
\end{proof}

For the remainder of this section, let $s = [K:\Q].$ We proceed to prove Theorem \ref{lowerbound2}.


\begin{proof}[Proof of Theorem \ref{lowerbound2}]
    We partition $I_K$ in two steps, as follows. Define the map $C_2 : I_K \rightarrow \{0,1\}$ as follows. For any ideal $I \in I_K,$ if there exists a prime ideal $\p \mid I$ such that $N(\p) \leq s-1,$ then $C_2(I) = 0.$ Otherwise, $C_2(I) = 1.$
    In addition, define the map $C_1: I_K \rightarrow \mathcal{P}(\N)$ by 
    \begin{align*}
        C_1(I) = \{N(\mathfrak{d}) : \dd \in I_K, \mathfrak{d} \mid I, \text{ and } s \leq N(\mathfrak{d}) \leq \sqrt{r-2}\}.
    \end{align*}
    Finally, define 
    $$A(I) = \begin{cases}
        \{0\} &\text{if }C_2(I) = 0\\
        C_1(I) &\text{if }C_2(I) = 1
    \end{cases}.$$ Note that this determines a partition of $I_K$ whose parts are $A^{-1}(S)$ for $S \in A(I_K).$ We wish to show that for any distinct $S_1, S_2 \in A(I_K),$ the sets $\sigma_{r,K}(A^{-1}(S_1))$ and $\sigma_{r,K}(A^{-1}(S_2))$ have disjoint closures. Write $S_1 = A(\fa)$ and $S_2 = A(\bb)$ for distinct $\fa,\bb \in I_K.$ Suppose $C_2(\fa)$ and $C_2(\bb)$ are not both 1; in particular, suppose without loss of generality that $C_2(\fa) = 0$ and $C_2(\bb) = 1.$ Then for any $\fa'$ such that $A(\fa) = A(\fa')$, we have $\sigma_{r,K}(\fa') \geq 1 + (s-1)^{-r}.$ 
    
    Meanwhile, we find an upper bound for $\sigma_{r,K}(\bb'),$ where $\bb'$ is any ideal such that $A(\bb) = A(\bb').$ Let $\dd$ be a divisor of $\bb'$ that is not the trivial ideal $\mathcal{O}_K$.
    We can factor $N(\dd) = \prod_{i=1}^k p_i^{e_i},$ where $p_i$ are distinct rational primes and $e_i \geq 1$. Then the number of distinct ideals with norm $N(\dd)$ is at most $\prod_{i=1}^k s^{e_i},$ because there are at most $s$ distinct prime ideals lying over any rational prime. Now, since $C_2(\bb') = 1,$ no integer in the interval $[2,s-1]$ divides $N(\bb')$ and therefore $N(\dd)$. Hence, each $p_i \geq s,$ so $ N(\dd) \geq \prod_{i=1}^k s^{e_i}.$ 

    Then we see that $$\sigma_{r,K}(\bb') \leq 1 + \sum_{d=s}^\infty d \cdot d^{-r} < 1 + \frac{1}{r-2} (s-1)^{-r+2}$$ by Lemma \ref{ideal1.75}. Notice that if $r > 2 + (s-1)^2,$ we obtain $$1 + (s-1)^{-r} > 1 + \frac{1}{r-2} (s-1)^{-r+2},$$ meaning $\sigma_{r,K}(A^{-1}(A(\fa)))$ and $\sigma_{r,K}(A^{-1}(A(\bb)))$ have disjoint closures when $C_2(\fa)$ and $C_2(\bb)$ are not both 1. 

    Now, suppose $C_2(\fa)  = C_2(\bb) = 1.$ Then $C_1(\fa) \neq C_1(\bb).$ Furthermore, suppose without loss of generality that the least-norm ideal $\dd$ dividing precisely one of $\fa$ and $\bb$ and satisfying $N(\dd) \leq \sqrt{r-2}$ divides $\fa.$ Let $d_0 = N(\dd).$ Note that $d_0 \geq s$. For any $\fa'$ such that $A(\fa) = A(\fa')$ and any $\bb'$ such that $A(\bb) = A(\bb'),$ we have $$\sigma_{r,K}(\fa') - \sigma_{r,K}(\bb') \geq d_0^{-r} - \sum_{d = d_0+1}^\infty d \cdot d^{-r}.$$ Again applying Lemma \ref{ideal1.75} shows that when $d_0 \leq \sqrt{r-2},$ this expression is positive, so $\sigma_{r,K}(A^{-1}(A(\fa)))$ and $\sigma_{r,K}(A^{-1}(A(\bb)))$ have disjoint closures when $C_2(\fa)=C_2(\bb)=1.$

    In conclusion, $\overline{\sigma_{r,K}(I_K)}$ has at least $A(I_K)$ many connected components. 
    We seek a lower bound on the size of $A(I_K).$ For any rational prime $s^{\frac{1}{s}} \leq p\leq (r-2)^{\frac{1}{2s}}$, take any prime ideal $\mathfrak{p}$ of $\mathcal{O}_K$ lying over $p\Z.$ Note that then $N(\mathfrak{p}) \leq \sqrt{r-2}$. So any subset of rational primes in the interval $[s^{\frac{1}{s}}, (r-2)^{\frac{1}{2s}}]$ corresponds uniquely to an element of $A(I_K).$ The lower bound follows.
\end{proof}

\begin{corollary}
    If $s$ is prime, then as $r$ approaches infinity, $C_{r,K} = \Omega\left( 2^{\frac{1}{s^2} \pi(\sqrt{r-2})}\right)$.
\end{corollary}

\begin{proof}
    Let $m = \sqrt{r-2}.$ Let $S_m$ be the set of prime ideals $\mathfrak{p}$ such that $N(\mathfrak{p}) \leq m,$ and let $U_m$ be the set of ideals in $S_m$ that have ramification index 1 and inertial degree 1. Let $N(S_m)$ and $N(U_m)$ be the set of norms of ideals in $S_m$ and $U_m$, respectively. By Chebotarev's density theorem, as $r$ approaches infinity, the proportion of elements in $S_m$ that are in $U_m$ is $\frac{1}{s}.$ Because $s$ is prime, each prime ideal in $U_m$ has exactly $s$ conjugate prime ideals, and except for finitely many ramified primes, each prime ideal in $S_m \setminus U_m$ has exactly 1 conjugate prime ideal (being itself). As a result, the proportion of elements in $N(S_m)$ that lie in $N(U_m)$ is $\frac{1}{s^2}.$ Finally, $N(S_m) \leq \pi(m)$ because the norm of a prime ideal $\mathfrak{p}$ is always at least the rational prime that $\mathfrak{p}$ lies over in $\Z$, so the proportion of rational primes less than or equal to $m$ that lie in $N(U_m)$ is at least $\frac{1}{s^2}.$ Each subset of $N(U_m)$ lies in $A(I_K),$ so the conclusion follows.
\end{proof}

\end{comment}
\section{Acknowledgments}

This research was conducted at the Duluth REU at the University of Minnesota Duluth, which is supported by National Science Foundation grant No.~DMS-2409861, Jane
Street Capital, and donations from Ray Sidney and Eric Wepsic. I would like
to thank Colin Defant, Noah Kravitz, Mitchell Lee, Maya Sankar, and Carl Schildkraut for providing guidance during the research process, as well as Derek Liu for helpful discussions. I would also like to thank Jonas Iskander, Noah Kravitz, Andrew Kwon, Mitchell Lee, and Yelena Mandelshtam for their detailed feedback throughout the editing process. Finally, I thank Colin
Defant and Joe Gallian for providing this wonderful opportunity to participate in the Duluth REU. 

	
\begin{thebibliography}{20}
        \bibitem{AB19} N.~Achenjang, A.~Berger: \emph{On gaps in the closures of images of divisor functions}, International Journal of Number Theory \textbf{15} (2019) 1023--1036.
        
        \bibitem{D15} C.~Defant: \textit{On the density of ranges of generalized divisor functions,} Notes on Number Theory and Discrete Mathematics \textbf{21} (2015), 87--88.

        \bibitem{D17} C.~Defant: \textit{Connected components of complex divisor functions,} Journal of Number Theory \textbf{190} (2018), 56--71.

        \bibitem{D16} C. Defant: \textit{Complex divisor functions}, Anal. Geom. Number Theory \textbf{1} (2016), 22--48.

        \bibitem{L86} R.~Laatsch: \textit{Measuring the abundancy of integers}, Mathematics Magazine \textbf{59(2)} (1986), 84--92.
		
		\bibitem{S17} C.~Sanna: \emph{On the closure of the image of the generalized divisor function}, Uniform Distribution Theory \textbf{12(2)} (2017) 77--90.

        

        \bibitem{W00} P.~A.~Weiner: \textit{The abundancy ratio, a measure of perfection,} Mathematics Magazine \textbf{73(4)} (2000), 307--310.

        \bibitem{Z18} N.~Zubrilina: \textit{On the number of connected components of ranges of divisor functions,} arXiv:1711.02871 (2017).
	\end{thebibliography}


\end{document}
